\chapter[Cox Coordinates and Quotient Constructions]
{Cox Coordinates and Quotient\\ Constructions}
\label{chap:cox-coordinates-quotients}

\begin{chapteroverview}
Projective space is built from homogeneous coordinates by deleting the
origin and quotienting by one scalar action:
\[
\PP^n
=
\bigl(\AAA^{n+1}\setminus\{0\}\bigr)/\Gm.
\]
A toric variety admits a striking generalization of this construction.
Every ray \(\rho\in\Sigma(1)\) contributes a coordinate \(x_\rho\), and
the polynomial ring
\[
S_\Sigma=\Bbbk[x_\rho:\rho\in\Sigma(1)]
\]
is graded by the divisor class group:
\[
\deg(x_\rho)=[D_\rho]\in\Cl(X_\Sigma).
\]
The grading records the independent rescalings of the coordinates.
After deleting a fan-determined exceptional set and quotienting by the
corresponding diagonal group, we recover \(X_\Sigma\).

The chapter builds toward
\[
\boxed{
X_\Sigma
\simeq
\bigl(\AAA^{\Sigma(1)}\setminus Z(\Sigma)\bigr)/G,
\qquad
G=\Hom(\Cl(X_\Sigma),\Gm).
}
\]
The quotient must be interpreted with care when the fan is not
simplicial; we will make the distinction precise when it arises.
\end{chapteroverview}

\begin{threedictionaries}
Cox coordinates package the fan, the boundary, and a multigraded
polynomial ring into one construction.

\medskip

\begin{center}
\small
\begin{tabularx}{\linewidth}{@{}
  >{\raggedright\arraybackslash}X
  >{\raggedright\arraybackslash}X
  >{\raggedright\arraybackslash}X@{}}
\toprule
\textbf{Algebra}
&
\textbf{Geometry}
&
\textbf{Combinatorics}
\\
\midrule
variable \(x_\rho\)
&
invariant divisor \(D_\rho\)
&
ray \(\rho\)
\\
\(\Cl(X_\Sigma)\)-grading
&
independent coordinate rescalings
&
relations among ray generators
\\
irrelevant ideal
&
exceptional set to delete
&
collections of rays missing from cones
\\
homogeneous quotient
&
toric variety \(X_\Sigma\)
&
the fan \(\Sigma\)
\\
\bottomrule
\end{tabularx}
\end{center}
\end{threedictionaries}

\begin{surfacelaboratory}
The surface examples will keep the construction concrete:
\[
\PP^2,
\qquad
\PP^1\times\PP^1,
\qquad
\mathbb F_r,
\qquad
\Blp{\PP^2}{p}.
\]
Their Cox rings are polynomial rings with three or four variables, but
the gradings and deleted coordinate loci differ. Those differences
encode the divisor classes and the fan.
\end{surfacelaboratory}

% Section 6.1 is inlined below so this chapter file is self-contained.
\section{Why homogeneous coordinates need a grading}
\label{sec:why-homogeneous-coordinates-need-grading}
\index{homogeneous coordinates!toric}
\index{class-group grading}
\index{Cox coordinate}

The last chapter ended with a polynomial ring waiting in the wings:
\[
S=\Bbbk[x_\rho:\rho\in\Sigma(1)].
\]
There is one variable for every ray, hence one variable for every
irreducible component \(D_\rho\) of the toric boundary. At first sight,
this looks almost too easy. A polynomial ring is familiar; why should
it contain the geometry of an arbitrary toric variety?

The subtlety is not the ring itself. The subtlety is the word
\emph{homogeneous}. For projective space, all coordinates may be
rescaled by one common scalar, and ordinary degree records that one
scaling direction. A general toric variety can have several
independent scaling directions. Its homogeneous coordinates therefore
need a grading by a group larger than \(\ZZ\).

\begin{agcheatsheet}
A point of projective space is an equivalence class
\[
[z_0:\cdots:z_n]
=
[\lambda z_0:\cdots:\lambda z_n],
\qquad \lambda\in\Bbbk^\times.
\]
Equivalently,
\[
\PP^n
=
\bigl(\AAA^{n+1}\setminus\{0\}\bigr)/\Gm,
\]
where \(\Gm\) acts by simultaneous scalar multiplication.

A polynomial equation
\[
f(z_0,\dots,z_n)=0
\]
is well defined on projective points only when replacing
\((z_0,\dots,z_n)\) by \((\lambda z_0,\dots,\lambda z_n)\) multiplies
\(f\) by one common power of \(\lambda\). This is exactly what happens
when every monomial of \(f\) has the same total degree:
\[
f(\lambda z_0,\dots,\lambda z_n)
=
\lambda^d f(z_0,\dots,z_n).
\]
Thus ordinary homogeneity is the algebraic record of one scalar
symmetry.
\end{agcheatsheet}

\subsection{One scaling direction: projective space}

For \(\PP^2\), write homogeneous coordinates as
\[
[x_0:x_1:x_2].
\]
The triples
\[
(x_0,x_1,x_2)
\qquad\text{and}\qquad
(\lambda x_0,\lambda x_1,\lambda x_2)
\]
represent the same point. Therefore the three variables have the same
weight:
\[
\deg(x_0)=\deg(x_1)=\deg(x_2)=1.
\]
A polynomial such as
\[
x_0x_1+x_2^2
\]
is homogeneous of degree \(2\), since both terms acquire the factor
\(\lambda^2\). By contrast,
\[
x_0+x_1x_2
\]
is not homogeneous: its two terms acquire different powers of
\(\lambda\), so its zero set is not preserved by rescaling.

This familiar story can be read in divisor language. The three
coordinate hyperplanes define linearly equivalent divisors:
\[
[D_0]=[D_1]=[D_2]=H\in\Cl(\PP^2).
\]
Thus the assignment
\[
\deg(x_i)=[D_i]
\]
recovers the ordinary grading after identifying \(\Cl(\PP^2)\simeq\ZZ\)
with generator \(H\).

\begin{pauseandthink}
Explain directly from the fan of \(\PP^2\) why the three invariant
boundary divisors have the same class. Which principal divisors give
the two independent relations?
\end{pauseandthink}

\subsection{Two scaling directions: \texorpdfstring{\(\PP^1\times\PP^1\)}{P1 x P1}}

Now consider
\[
\PP^1\times\PP^1.
\]
A point is written
\[
([x_0:x_1],[y_0:y_1]).
\]
There are two independent rescalings:
\[
(x_0,x_1;y_0,y_1)
\sim
(\lambda x_0,\lambda x_1;\mu y_0,\mu y_1),
\qquad
(\lambda,\mu)\in(\Bbbk^\times)^2.
\]
The \(x\)-coordinates respond only to \(\lambda\), while the
\(y\)-coordinates respond only to \(\mu\). One integer can no longer
record both behaviors. The natural grading is the bidegree
\[
\deg(x_0)=\deg(x_1)=(1,0),
\qquad
\deg(y_0)=\deg(y_1)=(0,1).
\]

\begin{figure}[ht]
\centering
\begin{tikzpicture}[scale=0.95]
  % P2 panel
  \begin{scope}[xshift=-4.4cm]
    \node[font=\bfseries] at (0,2.65) {one scaling direction};
    \node[font=\bfseries] at (0,2.18) {\(\PP^2\)};

    \draw[rounded corners,fill=DictionaryBlue!8,draw=DictionaryBlue,very thick]
      (-2.15,-0.15) rectangle (2.15,1.35);
    \node at (0,0.87) {\((x_0,x_1,x_2)\)};
    \node[DictionaryBlue,font=\large] at (0,0.28)
      {\(\lambda\cdot(x_0,x_1,x_2)\)};

    \draw[->,very thick,DictionaryBlue]
      (-2.55,0.6)--(-2.18,0.6);
    \node[DictionaryBlue,left] at (-2.55,0.6) {\(\lambda\)};

    \node[align=center] at (0,-0.72)
      {all variables have\\the same degree};
    \node[fill=DictionaryBlue!10,rounded corners,inner sep=3pt]
      at (0,-1.58) {grading group \(\ZZ\)};
  \end{scope}

  % P1xP1 panel
  \begin{scope}[xshift=4.4cm]
    \node[font=\bfseries] at (0,2.65) {two scaling directions};
    \node[font=\bfseries] at (0,2.18) {\(\PP^1\times\PP^1\)};

    \draw[rounded corners,fill=PredictionTeal!8,draw=PredictionTeal,very thick]
      (-2.35,0.62) rectangle (-0.15,1.35);
    \draw[rounded corners,fill=ExampleGold!12,draw=ExampleGold,very thick]
      (0.15,0.62) rectangle (2.35,1.35);
    \node at (-1.25,0.98) {\((x_0,x_1)\)};
    \node at (1.25,0.98) {\((y_0,y_1)\)};

    \draw[->,very thick,PredictionTeal]
      (-2.75,0.98)--(-2.38,0.98);
    \node[PredictionTeal,left] at (-2.75,0.98) {\(\lambda\)};
    \draw[->,very thick,ExampleGold]
      (2.75,0.98)--(2.38,0.98);
    \node[ExampleGold,right] at (2.75,0.98) {\(\mu\)};

    \node[PredictionTeal] at (-1.25,0.18) {degree \((1,0)\)};
    \node[ExampleGold] at (1.25,0.18) {degree \((0,1)\)};
    \node[align=center] at (0,-0.72)
      {the two coordinate blocks\\scale independently};
    \node[fill=PredictionTeal!8,rounded corners,inner sep=3pt]
      at (0,-1.58) {grading group \(\ZZ^2\)};
  \end{scope}
\end{tikzpicture}
\caption{A grading records how coordinate blocks transform under
rescaling. Projective space has one common scaling direction;
\(\PP^1\times\PP^1\) has two independent directions.}
\label{fig:one-scaling-versus-two}
\end{figure}

The polynomial
\[
x_0y_0+x_1y_1
\]
has bidegree \((1,1)\). Under rescaling, both terms acquire the same
factor:
\[
(\lambda x_0)(\mu y_0)+(\lambda x_1)(\mu y_1)
=
\lambda\mu(x_0y_0+x_1y_1).
\]
Its zero set is therefore well defined on \(\PP^1\times\PP^1\).

By contrast, an ordinary total-degree grading would declare
\[
x_0+y_0
\]
to be homogeneous of degree \(1\). But
\[
x_0+y_0
\longmapsto
\lambda x_0+\mu y_0,
\]
and there is no common factor unless \(\lambda=\mu\). The equation
\(x_0+y_0=0\) is not well defined on the product. Ordinary degree has
forgotten which scaling acts on which coordinate.

\begin{pauseandthink}
Determine which of the following polynomials are bihomogeneous, and
record their bidegrees when they are:
\[
x_0^2y_0+x_0x_1y_1,
\qquad
x_0y_0+x_1y_1^2,
\qquad
x_0x_1+y_0y_1.
\]
For each failure, identify the two monomials whose scaling behaviors
do not match.
\end{pauseandthink}

The divisor class group explains the bidegree. If \(F_1\) and \(F_2\)
are the two ruling classes, then
\[
\Cl(\PP^1\times\PP^1)
\simeq
\ZZ F_1\oplus\ZZ F_2,
\]
and the four invariant boundary divisors satisfy
\[
[D_{x_0}]=[D_{x_1}]=F_1,
\qquad
[D_{y_0}]=[D_{y_1}]=F_2.
\]
Hence the bidegree is again obtained by the rule
\[
\deg(x_\rho)=[D_\rho].
\]

\subsection{One variable for every boundary divisor}

The examples suggest the general construction. Let \(X_\Sigma\) be a
toric variety whose rays span \(N_\RR\); in particular, this holds for
every complete toric variety. For each ray
\[
\rho\in\Sigma(1),
\]
introduce a variable \(x_\rho\). We want the equation
\[
x_\rho=0
\]
to remember the invariant prime divisor \(D_\rho\). More generally, a
monomial
\[
x^a=\prod_{\rho\in\Sigma(1)}x_\rho^{a_\rho}
\]
records the effective invariant divisor
\[
D_a=\sum_{\rho\in\Sigma(1)}a_\rho D_\rho.
\]

This makes the correct notion of degree unavoidable:
\[
\boxed{
\deg(x^a)
=
[D_a]
=
\left[\sum_{\rho}a_\rho D_\rho\right]
\in\Cl(X_\Sigma).
}
\]
In particular,
\[
\boxed{
\deg(x_\rho)=[D_\rho].
}
\]

The class group is not being added as decoration. It is precisely the
group that remembers when two invariant divisors differ by a principal
divisor. By \cref{thm:toric-class-group-sequence}, we have the exact
sequence
\[
0\longrightarrow M
\longrightarrow \ZZ^{\Sigma(1)}
\stackrel{\deg}{\longrightarrow}\Cl(X_\Sigma)
\longrightarrow0,
\]
where
\[
m\longmapsto
\bigl(\langle m,u_\rho\rangle\bigr)_\rho
\]
and
\[
(a_\rho)_\rho
\longmapsto
\left[\sum_\rho a_\rho D_\rho\right].
\]
The first map records the principal-divisor relation
\[
\divchi{m}
=
\sum_\rho\langle m,u_\rho\rangle D_\rho.
\]
Thus two monomials \(x^a\) and \(x^b\) have the same degree exactly
when
\[
\sum_\rho(a_\rho-b_\rho)D_\rho
\]
is principal.

\begin{figure}[ht]
\centering
\begin{tikzpicture}[node distance=1.5cm]
  \node[draw=DictionaryBlue,very thick,rounded corners,
        fill=DictionaryBlue!8,minimum width=1.65cm,minimum height=1cm]
    (ray) {ray \(\rho\)};
  \node[draw=PredictionTeal,very thick,rounded corners,
        fill=PredictionTeal!8,minimum width=1.95cm,minimum height=1cm,
        right=of ray]
    (var) {variable \(x_\rho\)};
  \node[draw=ExampleGold,very thick,rounded corners,
        fill=ExampleGold!12,minimum width=1.95cm,minimum height=1cm,
        right=of var]
    (div) {divisor \(D_\rho\)};
  \node[draw=SurfaceGreen,very thick,rounded corners,
        fill=SurfaceGreen!10,minimum width=2.15cm,minimum height=1cm,
        right=of div]
    (class) {degree \([D_\rho]\)};

  \draw[->,very thick] (ray)--(var)
    node[midway,above=7pt,font=\scriptsize,align=center]
      {name a\\coordinate};
  \draw[->,very thick] (var)--(div)
    node[midway,above=7pt,font=\scriptsize,align=center]
      {set it\\to zero};
  \draw[->,very thick] (div)--(class)
    node[midway,above=7pt,font=\scriptsize,align=center]
      {pass to the\\divisor class};

  \node[align=center,below=0.95cm of var,xshift=2.0cm]
    {the fan supplies the variables;\\the class group supplies their degrees};
\end{tikzpicture}
\caption{The basic chain behind Cox coordinates. Each ray gives both a
boundary divisor and a variable; the divisor class becomes the
variable's degree.}
\label{fig:ray-variable-divisor-degree}
\end{figure}

\begin{definition}[The class-group grading]
\label{def:class-group-grading-preview}
Let
\[
S_\Sigma=\Bbbk[x_\rho:\rho\in\Sigma(1)].
\]
Give \(S_\Sigma\) the \(\Cl(X_\Sigma)\)-grading determined by
\[
\deg(x_\rho)=[D_\rho].
\]
A polynomial is \emph{homogeneous of degree}
\(\beta\in\Cl(X_\Sigma)\) if every monomial appearing in it has degree
\(\beta\).
\end{definition}

\noindent
\textit{Reference.}
This is the grading introduced in Cox--Little--Schenck, \S5.2,
``The Grading.'' Their exact sequence~(5.1.1) is the same class-group
sequence used above.

\begin{observation}
There is optional representation-theoretic language for this picture.
If an abelian group \(A\) grades a ring
\[
S=\bigoplus_{\alpha\in A}S_\alpha,
\]
then the group
\[
\Hom(A,\Gm)
\]
acts diagonally on \(\Spec(S)\), and the spaces \(S_\alpha\) are its
weight spaces. In our setting,
\[
A=\Cl(X_\Sigma).
\]
Thus the class-group grading is the algebraic shadow of the rescalings
that will eventually be quotiented out. No familiarity with weight
spaces is needed: the concrete content is simply that all monomials of
a homogeneous polynomial acquire the same common factor.
\end{observation}

\subsection{Surface laboratory: the blowup grading}

Return to the toric blowup
\[
X=\Blp{\PP^2}{p}
\]
from Chapter~5. Its rays are generated by
\[
u_0=-e_1-e_2,
\qquad
u_1=e_1,
\qquad
u_E=e_1+e_2,
\qquad
u_2=e_2.
\]
Let the corresponding variables be
\[
x_0,x_1,x_E,x_2.
\]
The principal-divisor relations are
\[
-D_0+D_1+E\sim0,
\qquad
-D_0+E+D_2\sim0.
\]
Writing \(H=[D_0]\), we obtain
\[
[D_1]=[D_2]=H-E.
\]
Therefore the polynomial ring
\[
S_X=\Bbbk[x_0,x_1,x_E,x_2]
\]
has the \(\Cl(X)\simeq\ZZ H\oplus\ZZ E\) grading
\[
\deg(x_0)=H,
\qquad
\deg(x_1)=\deg(x_2)=H-E,
\qquad
\deg(x_E)=E.
\]

Now consider the divisor class \(2H-E\). The five monomials
\[
x_0x_1,
\qquad
x_0x_2,
\qquad
x_1^2x_E,
\qquad
x_1x_2x_E,
\qquad
x_2^2x_E
\]
all have degree \(2H-E\). They are the Cox-coordinate counterparts of
the five lattice points in the truncated polygon \(P_{2H-E}\) from
Chapter~5.

This is our first glimpse of what the Cox ring will accomplish: one
multigraded polynomial ring simultaneously contains the monomial
sections of many divisor classes.

\begin{prediction}
For the divisor class \(dH-rE\), predict the exponents of a general
monomial
\[
x_0^{a_0}x_1^{a_1}x_E^{a_E}x_2^{a_2}
\]
of degree \(dH-rE\). How should the resulting integer conditions be
related to the inequalities
\[
a,b\geq0,
\qquad
r\leq a+b\leq d
\]
that define the section polytope \(P_{dH-rE}\)?
\end{prediction}

\begin{threedictionaries}
The need for a class-group grading can now be read in all three
languages.

\medskip

\begin{center}
\small
\begin{tabularx}{\linewidth}{@{}
  >{\raggedright\arraybackslash}X
  >{\raggedright\arraybackslash}X
  >{\raggedright\arraybackslash}X@{}}
\toprule
\textbf{Algebra}
&
\textbf{Geometry}
&
\textbf{Combinatorics}
\\
\midrule
variable \(x_\rho\)
&
boundary divisor \(D_\rho\)
&
ray \(\rho\)
\\
monomial \(x^a\)
&
effective invariant divisor \(\sum a_\rho D_\rho\)
&
nonnegative exponent vector \(a\in\ZZ^{\Sigma(1)}\)
\\
degree in \(\Cl(X_\Sigma)\)
&
linear-equivalence class
&
exponent vectors modulo principal-divisor relations
\\
homogeneous polynomial
&
a scaling-invariant equation
&
monomials lying in one fiber of the degree map
\\
\bottomrule
\end{tabularx}
\end{center}
\end{threedictionaries}

The lesson is compact:
\[
\boxed{
\begin{gathered}
\text{one scaling direction gives ordinary degree;}\\
\text{many scaling directions give the class-group grading.}
\end{gathered}
}
\]

\begin{guidingquestion}
For a fixed divisor class \(\beta\in\Cl(X_\Sigma)\), what geometric
object is represented by the homogeneous piece
\[
(S_\Sigma)_\beta?
\]
Can it be identified with the global sections of a divisor in the
class \(\beta\), and how does multiplication in the polynomial ring
encode multiplication of sections?
\end{guidingquestion}
\section{The Cox ring and its homogeneous pieces}
\label{sec:cox-ring-homogeneous-pieces}
\index{Cox ring}
\index{homogeneous piece!of the Cox ring}
\index{global section!inside the Cox ring}

Section~6.1 explained why the variables of a toric homogeneous
coordinate ring must be graded by \(\Cl(X_\Sigma)\). We now take the
next step: instead of looking at one degree at a time, we place all
divisor classes and all their monomial sections inside a single
polynomial ring.

Throughout this section, the rays of \(\Sigma\) span \(N_\RR\). This
holds in particular when \(X_\Sigma\) is complete. Equivalently,
\(X_\Sigma\) has no torus factor, and the map
\[
M\longrightarrow \ZZ^{\Sigma(1)},
\qquad
m\longmapsto
\bigl(\langle m,u_\rho\rangle\bigr)_\rho
\]
is injective.

\subsection{One polynomial ring for all divisor classes}

\begin{definition}[The Cox ring]
\label{def:toric-cox-ring}
The \emph{Cox ring}, or \emph{total coordinate ring}, of
\(X_\Sigma\) is the polynomial ring
\[
\operatorname{Cox}(X_\Sigma)
=
S_\Sigma
=
\Bbbk[x_\rho:\rho\in\Sigma(1)]
\]
with the \(\Cl(X_\Sigma)\)-grading determined by
\[
\deg(x_\rho)=[D_\rho].
\]
For \(\alpha\in\Cl(X_\Sigma)\), the homogeneous piece of degree
\(\alpha\) is denoted
\[
(S_\Sigma)_\alpha.
\]
Thus
\[
S_\Sigma
=
\bigoplus_{\alpha\in\Cl(X_\Sigma)}
(S_\Sigma)_\alpha,
\qquad
(S_\Sigma)_\alpha(S_\Sigma)_\beta
\subseteq
(S_\Sigma)_{\alpha+\beta}.
\]
\end{definition}

\noindent
\textit{Reference.}
This is the total coordinate ring and grading of
Cox--Little--Schenck, \S5.2.

The ring is polynomial, but its grading is geometric. A monomial
\[
x^a=\prod_\rho x_\rho^{a_\rho}
\]
represents the effective invariant divisor
\[
E_a=\sum_\rho a_\rho D_\rho,
\]
and its degree is precisely the divisor class
\[
\deg(x^a)=[E_a].
\]
Consequently,
\[
(S_\Sigma)_\alpha
\]
is spanned by the monomials whose effective invariant divisors lie in
the class \(\alpha\).

\begin{example}[The familiar graded pieces of \(\PP^2\)]
For \(\PP^2\),
\[
S=\Bbbk[x_0,x_1,x_2],
\qquad
\deg(x_i)=H.
\]
After identifying \(\Cl(\PP^2)\simeq\ZZ H\), the homogeneous piece
\(S_{dH}\) is the usual vector space of degree-\(d\) homogeneous
polynomials. For example,
\[
S_{2H}
=
\Span_\Bbbk
\{x_0^2,x_0x_1,x_0x_2,x_1^2,x_1x_2,x_2^2\}.
\]
Thus ordinary homogeneous coordinates are the rank-one instance of
the Cox construction.
\end{example}

\begin{example}[Bihomogeneous pieces on \(\PP^1\times\PP^1\)]
For
\[
S=\Bbbk[x_0,x_1,y_0,y_1]
\]
with
\[
\deg(x_i)=(1,0),
\qquad
\deg(y_j)=(0,1),
\]
the piece \(S_{(a,b)}\) is spanned by the bihomogeneous monomials
\[
x_0^i x_1^{a-i}y_0^j y_1^{b-j},
\qquad
0\leq i\leq a,
\quad
0\leq j\leq b.
\]
Hence
\[
\dim_\Bbbk S_{(a,b)}=(a+1)(b+1)
\]
when \(a,b\geq0\), while the piece is zero if either coordinate is
negative.
\end{example}

These examples suggest that a homogeneous piece should be more than an
abstract collection of polynomials: it should be the space of sections
of the corresponding divisor class.

\subsection{A degree fiber is a section polytope}

Choose a torus-invariant Weil divisor
\[
D=\sum_{\rho\in\Sigma(1)}d_\rho D_\rho,
\qquad
\alpha=[D]\in\Cl(X_\Sigma).
\]
Recall from Chapter~5 that
\[
P_D
=
\left\{
 m\in M_\RR:
 \langle m,u_\rho\rangle+d_\rho\geq0
 \text{ for every }\rho
\right\}.
\]
For a lattice point \(m\in P_D\cap M\), define the monomial
\[
x^{\langle m,D\rangle}
:=
\prod_{\rho\in\Sigma(1)}
 x_\rho^{\langle m,u_\rho\rangle+d_\rho}.
\]
The inequalities defining \(P_D\) guarantee that every exponent is
nonnegative.

\begin{lemma}[Homogenizing a character]
\label{lem:homogenizing-character}
For every \(m\in P_D\cap M\), the monomial
\(x^{\langle m,D\rangle}\) has degree \([D]\).
\end{lemma}

\begin{proof}
Its degree is
\[
\left[
\sum_\rho
\bigl(\langle m,u_\rho\rangle+d_\rho\bigr)D_\rho
\right]
=
\left[D+\divchi{m}\right]
=[D].
\]
The principal-divisor formula is the entire reason the degree does not
change when the character \(\character m\) is homogenized.
\end{proof}

The exponent vector of this monomial is
\[
\bigl(\langle m,u_\rho\rangle+d_\rho\bigr)_\rho
\in\ZZ_{\geq0}^{\Sigma(1)}.
\]
Thus the lattice points of \(P_D\) are converted into the nonnegative
integer points lying over \([D]\) under the degree map.

\begin{figure}[ht]
\centering
\begin{tikzpicture}[scale=0.92]
  % polytope side
  \begin{scope}[xshift=-4.15cm]
    \draw[->,black!60] (-0.25,0)--(2.55,0) node[below] {\(a\)};
    \draw[->,black!60] (0,-0.25)--(0,2.55) node[left] {\(b\)};
    \filldraw[fill=DictionaryBlue!12,draw=DictionaryBlue,very thick]
      (1,0)--(2,0)--(0,2)--(0,1)--cycle;
    \foreach \p in {(1,0),(2,0),(1,1),(0,2),(0,1)} {
      \fill \p circle (2.3pt);
    }
    \node[font=\bfseries] at (1,3.0) {lattice points of \(P_{2H-E}\)};
    \node[align=center,font=\small] at (1,-0.7)
      {five points in a\\truncated triangle};
  \end{scope}

  % arrow
  \draw[->,very thick,PredictionTeal]
    (-1.35,1.22)--(1.25,1.22);
  \node[PredictionTeal,fill=white,inner sep=2pt]
    at (-0.05,1.62) {homogenize};
  \node[font=\scriptsize,align=center]
    at (-0.05,0.70)
    {the same\\integer data};

  % degree fiber side
  \begin{scope}[xshift=4.15cm]
    \draw[rounded corners,fill=PredictionTeal!8,
          draw=PredictionTeal,very thick]
      (-2.35,-0.18) rectangle (2.35,2.42);
    \node[font=\bfseries] at (0,2.78)
      {monomials of degree \(2H-E\)};
    \node[align=left,font=\small] at (0,1.15) {$
      \begin{aligned}
      x_0x_1,\qquad &x_1^2x_E,\\
      x_1x_2x_E,\qquad &x_2^2x_E,\\
      &x_0x_2
      \end{aligned}$};
    \node[align=center,font=\small] at (0,-0.62)
      {five nonnegative exponent vectors\\in one fiber of the degree map};
  \end{scope}
\end{tikzpicture}
\caption{For \(D=2H-E\) on \(\Blp{\PP^2}{p}\), homogenization turns
the five lattice points of \(P_D\) into the five monomials in the
homogeneous piece \(S_{2H-E}\).}
\label{fig:polytope-to-degree-fiber}
\end{figure}

\subsection{Homogeneous pieces are spaces of sections}

\begin{theorem}[Sections inside the Cox ring]
\label{thm:cox-piece-global-sections}
Let \(D\) be a torus-invariant Weil divisor on a toric variety without
torus factors, and set \(\alpha=[D]\). Then the assignment
\[
\character m
\longmapsto
x^{\langle m,D\rangle}
\]
induces an isomorphism
\[
\boxed{
H^0(X_\Sigma,\mathcal O_{X_\Sigma}(D))
\simeq
(S_\Sigma)_\alpha.
}
\]
Under this isomorphism, the character basis indexed by
\(P_D\cap M\) becomes the monomial basis of the homogeneous piece.
\end{theorem}

\noindent
\textit{Reference.}
This is Cox--Little--Schenck, Proposition~5.4.1.

\begin{proof}
By \cref{thm:lattice-points-global-sections}, the left-hand side has
basis
\[
\{\character m:m\in P_D\cap M\}.
\]
By \cref{lem:homogenizing-character}, each basis element maps to a
monomial of degree \(\alpha\).

Conversely, let
\[
x^a=\prod_\rho x_\rho^{a_\rho}
\]
be any monomial of degree \(\alpha\). Then
\[
\left[\sum_\rho a_\rho D_\rho\right]=[D],
\]
so the invariant divisor
\[
\sum_\rho a_\rho D_\rho-D
\]
is principal. Therefore it equals \(\divchi m\) for some \(m\in M\),
and hence
\[
a_\rho=\langle m,u_\rho\rangle+d_\rho
\qquad
\text{for every }\rho.
\]
Because all \(a_\rho\geq0\), we have \(m\in P_D\cap M\). The
assumption that the rays span \(N_\RR\) makes \(m\) unique. Thus the
two monomial bases are in bijection.
\end{proof}

This theorem answers the guiding question from the previous section:
\[
\boxed{
(S_\Sigma)_{[D]}
\text{ is a canonical polynomial model for }
H^0(X_\Sigma,\mathcal O(D)).
}
\]
The explicit isomorphism depends on the chosen representative \(D\),
but the homogeneous piece depends only on its class.

\begin{warningbox}
The notation \(S_{[D]}\) does not mean that every class has a
nonzero homogeneous piece. It is possible that
\[
(S_\Sigma)_\alpha=0.
\]
Geometrically, this means that the class \(\alpha\) contains no
effective torus-invariant divisor, equivalently that the corresponding
section polyhedron has no lattice points.
\end{warningbox}

\subsection{Multiplication of sections becomes multiplication of monomials}

Let \(D\) and \(E\) be torus-invariant divisors with classes
\[
\alpha=[D],
\qquad
\beta=[E].
\]
There is a natural multiplication map
\[
H^0(X_\Sigma,\mathcal O(D))
\otimes
H^0(X_\Sigma,\mathcal O(E))
\longrightarrow
H^0(X_\Sigma,\mathcal O(D+E)).
\]
On character sections,
\[
\character m\cdot\character{m'}
=
\character{m+m'}.
\]
The corresponding Cox monomials satisfy
\[
x^{\langle m,D\rangle}
 x^{\langle m',E\rangle}
=
x^{\langle m+m',D+E\rangle}.
\]
Therefore the diagram
\[
\begin{tikzcd}[column sep=large,row sep=large]
H^0(\mathcal O(D))\otimes H^0(\mathcal O(E))
  \arrow[r]
  \arrow[d,"\sim"']
&
H^0(\mathcal O(D+E))
  \arrow[d,"\sim"]
\\
S_\alpha\otimes S_\beta
  \arrow[r]
&
S_{\alpha+\beta}
\end{tikzcd}
\]
commutes, where the bottom arrow is ordinary polynomial
multiplication.

This is why the direct sum over all divisor classes is so useful: the
Cox ring does not merely list section spaces. It remembers how their
sections multiply.

\begin{observation}
For a general normal variety, a Cox ring is often introduced as a
direct sum of section spaces over divisor classes. In the toric case,
\cref{thm:cox-piece-global-sections} identifies that direct sum with
the exceptionally concrete polynomial ring
\[
\Bbbk[x_\rho:\rho\in\Sigma(1)].
\]
The absence of algebraic relations among the variables is a special
feature of toric geometry: each variable independently records one
boundary component.
\end{observation}

\subsection{Surface laboratory: solving the blowup prediction}

For
\[
X=\Blp{\PP^2}{p},
\]
Section~6.1 found
\[
\deg(x_0)=H,
\qquad
\deg(x_1)=\deg(x_2)=H-E,
\qquad
\deg(x_E)=E.
\]
Consider a monomial
\[
x_0^{a_0}x_1^{a_1}x_E^{a_E}x_2^{a_2}.
\]
Its degree is
\[
(a_0+a_1+a_2)H
+
(a_E-a_1-a_2)E.
\]
It has degree \(dH-rE\) precisely when
\[
a_0+a_1+a_2=d,
\qquad
 a_E-a_1-a_2=-r.
\]
Set
\[
a=a_1,
\qquad
b=a_2.
\]
Then
\[
a_0=d-a-b,
\qquad
 a_E=a+b-r.
\]
The requirement that all four exponents be nonnegative becomes
\[
a,b\geq0,
\qquad
r\leq a+b\leq d.
\]
These are exactly the inequalities defining \(P_{dH-rE}\). Thus the
lattice point \((a,b)\) corresponds to the Cox monomial
\[
\boxed{
x_0^{d-a-b}x_1^a x_E^{a+b-r}x_2^b.
}
\]

The class-group grading and the section polytope are therefore two
ways of organizing the same integer solutions:
\[
\boxed{
P_{dH-rE}\cap M
\longleftrightarrow
\{\text{monomials of Cox degree }dH-rE\}.
}
\]

\begin{pauseandthink}
For \(D=3H-2E\), list the lattice points of \(P_D\) and the
corresponding Cox monomials. Verify the degree of each monomial
directly from the four variable degrees.
\end{pauseandthink}

\begin{threedictionaries}
A homogeneous piece of the Cox ring is the same object viewed in three
languages.

\medskip

\begin{center}
\small
\begin{tabularx}{\linewidth}{@{}
  >{\raggedright\arraybackslash}X
  >{\raggedright\arraybackslash}X
  >{\raggedright\arraybackslash}X@{}}
\toprule
\textbf{Algebra}
&
\textbf{Geometry}
&
\textbf{Combinatorics}
\\
\midrule
homogeneous piece \(S_{[D]}\)
&
sections of \(\mathcal O(D)\)
&
lattice points of \(P_D\)
\\
monomial \(x^{\langle m,D\rangle}\)
&
character section \(\character m\)
&
exponent vector over the degree fiber
\\
product in the Cox ring
&
product of sections
&
addition of lattice points and exponent vectors
\\
\bottomrule
\end{tabularx}
\end{center}
\end{threedictionaries}

The Cox ring has now collected every section space into one
multigraded polynomial ring. But this ring is the coordinate ring of
the full affine space
\[
\AAA^{\Sigma(1)}=\Spec(S_\Sigma),
\]
and that affine space is too large. Some coordinate patterns do not
represent any point of \(X_\Sigma\).

\begin{guidingquestion}
For projective space, the only forbidden coordinate pattern is
\((0,\dots,0)\). For a general fan, which collections of Cox
coordinates are forbidden to vanish simultaneously, and how can the
cones of \(\Sigma\) encode them as the zero set of an ideal?
\end{guidingquestion}
\section{The irrelevant ideal and the exceptional set}
\label{sec:irrelevant-ideal-exceptional-set}
\index{irrelevant ideal}
\index{exceptional set!Cox construction}
\index{primitive collection}
\index{Cox coordinate!vanishing test}

The Cox ring remembers the rays of the fan and the linear-equivalence
relations among their divisors. It does \emph{not} yet remember which
collections of rays span cones. That missing incidence data is exactly
what tells us which Cox coordinates are allowed to vanish together.

For projective space, the answer is familiar: all homogeneous
coordinates may vanish individually, but they may not vanish
simultaneously. Thus one deletes the origin before quotienting. For a
general toric variety, the forbidden locus can be a union of several
coordinate subspaces. The fan determines precisely which ones.

Throughout this section, write
\[
S_\Sigma=\Bbbk[x_\rho:\rho\in\Sigma(1)]
\]
and regard
\[
\AAA^{\Sigma(1)}=\Spec(S_\Sigma)
\]
as the affine space of Cox-coordinate tuples
\[
z=(z_\rho)_{\rho\in\Sigma(1)}.
\]

\subsection{The monomial attached to a cone}

Let \(\sigma\in\Sigma\). Define
\[
\boxed{
x^{\widehat\sigma}
=
\prod_{\rho\notin\sigma(1)}x_\rho.
}
\]
The hat means that we multiply the variables whose rays are
\emph{omitted} from \(\sigma\).

The principal open set
\[
D(x^{\widehat\sigma})
=
\{z:x^{\widehat\sigma}(z)\neq0\}
\]
consists of those coordinate tuples for which
\[
z_\rho\neq0
\qquad
\text{whenever }\rho\notin\sigma(1).
\]
The coordinates indexed by rays inside \(\sigma\) may vanish; all the
others must be invertible.

\begin{pauseandthink}
Suppose \(\sigma\) has rays \(\rho_1,\rho_2\) and the fan has five
rays in total. Write \(x^{\widehat\sigma}\). Which two coordinates may
vanish on \(D(x^{\widehat\sigma})\), and which three are forced to be
nonzero?
\end{pauseandthink}

If \(\tau\preceq\sigma\), then \(\tau(1)\subseteq\sigma(1)\), so
\[
x^{\widehat\tau}
\]
is divisible by \(x^{\widehat\sigma}\). Consequently, the monomials
coming from maximal cones already generate all the others.

\begin{definition}[The irrelevant ideal]
\label{def:irrelevant-ideal}
The \emph{irrelevant ideal} of the fan \(\Sigma\) is
\[
\boxed{
B(\Sigma)
=
\left\langle
x^{\widehat\sigma}:\sigma\in\Sigma
\right\rangle
=
\left\langle
x^{\widehat\sigma}:\sigma\in\Sigma_{\max}
\right\rangle
\subseteq S_\Sigma.
}
\]
Its zero set
\[
\boxed{
Z(\Sigma)=V(B(\Sigma))
\subseteq\AAA^{\Sigma(1)}
}
\]
is called the \emph{exceptional set}.
\end{definition}

\noindent
\textit{Reference.}
This is the construction in Cox--Little--Schenck, \S5.1, under
``The Exceptional Set.''

The name ``irrelevant'' comes from projective algebraic geometry. It
does not mean that the ideal is unimportant. Quite the opposite:
\(B(\Sigma)\) is the part of the Cox construction that remembers how
the rays are assembled into cones.

Since the complement of the zero set of an ideal generated by
\(f_1,\dots,f_s\) is the union of the principal opens \(D(f_i)\), we
obtain
\[
\boxed{
\AAA^{\Sigma(1)}\setminus Z(\Sigma)
=
\bigcup_{\sigma\in\Sigma_{\max}}
D(x^{\widehat\sigma}).
}
\]
Thus every maximal cone contributes one affine chart upstairs in Cox
coordinate space. In the next section, quotienting this chart will
recover the toric affine chart \(U_\sigma\).

\subsection{Which coordinates are allowed to vanish together?}

For a point \(z=(z_\rho)\in\AAA^{\Sigma(1)}\), define its
\emph{zero-coordinate set}
\[
I(z)=\{\rho\in\Sigma(1):z_\rho=0\}.
\]
The fan gives a remarkably simple test for whether \(z\) survives the
deletion of \(Z(\Sigma)\).

\begin{proposition}[The cone test for Cox coordinates]
\label{prop:cone-test-cox-coordinates}
A point \(z\in\AAA^{\Sigma(1)}\) lies outside \(Z(\Sigma)\) if and
only if there is a cone \(\sigma\in\Sigma\) such that
\[
\boxed{
I(z)\subseteq\sigma(1).
}
\]
In words: the coordinates that vanish must be indexed by rays that all
fit inside one cone of the fan.
\end{proposition}

\begin{proof}
The point \(z\) lies outside \(Z(\Sigma)\) exactly when
\[
x^{\widehat\sigma}(z)\neq0
\]
for at least one maximal cone \(\sigma\). This nonvanishing says that
\(z_\rho\neq0\) for every \(\rho\notin\sigma(1)\). Equivalently, every
ray whose coordinate vanishes belongs to \(\sigma(1)\), which is the
condition
\[
I(z)\subseteq\sigma(1).
\]
\end{proof}

This criterion makes the geometry of the deletion visible:
\[
\boxed{
\text{a collection of coordinates may vanish together}
\iff
\text{their rays lie together in a cone.}
}
\]

\begin{warningbox}
The grading alone does not determine this condition. The grading knows
the ray generators and their divisor-class relations. The irrelevant
ideal additionally knows \emph{which subsets of rays span cones}.
Thus two fans built from the same ray data can have the same Cox ring
and grading but different irrelevant ideals.
\end{warningbox}

\subsection{Minimal forbidden collections}

It is often inefficient to list every forbidden set of rays. Instead,
we record the minimal ones.

\begin{definition}[Primitive collection]
\label{def:primitive-collection}
A subset \(C\subseteq\Sigma(1)\) is a \emph{primitive collection} if

\begin{enumerate}[label=(\alph*),wide=0pt]
\item the rays in \(C\) are not all contained in any cone of
\(\Sigma\);

\item every proper subset of \(C\) is contained in the rays of some
cone of \(\Sigma\).
\end{enumerate}
\end{definition}

Thus a primitive collection is a minimal forbidden simultaneous
vanishing pattern.

\begin{proposition}[The exceptional set from primitive collections]
\label{prop:exceptional-set-primitive-collections}
The irreducible coordinate-subspace components of \(Z(\Sigma)\) are
indexed by the primitive collections:
\[
\boxed{
Z(\Sigma)
=
\bigcup_{C\text{ primitive}}
V(x_\rho:\rho\in C).
}
\]
\end{proposition}

\noindent
\textit{Reference.}
This is Cox--Little--Schenck, Definition~5.1.5 and
Proposition~5.1.6.

\begin{proof}[Conceptual proof sketch]
If all coordinates indexed by a primitive collection \(C\) vanish,
then their zero-coordinate set cannot lie in one cone. By
\cref{prop:cone-test-cox-coordinates}, every such point belongs to
\(Z(\Sigma)\).

Conversely, take a forbidden zero-coordinate set. Since the fan is
finite, it contains a minimal subset that still fails to lie in one
cone. That minimal subset is a primitive collection. Hence every point
of the exceptional set lies in one of the displayed coordinate
subspaces.
\end{proof}

\subsection{Projective space: delete only the origin}

For \(\PP^n\), the fan has rays
\[
\rho_0,\dots,\rho_n,
\]
and each maximal cone contains all but one of them. If \(\sigma_i\)
is the maximal cone omitting \(\rho_i\), then
\[
x^{\widehat\sigma_i}=x_i.
\]
Therefore
\[
B(\Sigma)=\langle x_0,\dots,x_n\rangle,
\qquad
Z(\Sigma)=\{0\}.
\]
The Cox open set is exactly
\[
\AAA^{n+1}\setminus\{0\},
\]
recovering the familiar starting point for homogeneous coordinates on
projective space.

The unique primitive collection is the full set
\[
\{\rho_0,\dots,\rho_n\}.
\]
Any proper subset spans a cone, but all \(n+1\) rays do not.

\subsection{Two independent projective factors}

For \(\PP^1\times\PP^1\), use Cox coordinates
\[
(x_0,x_1;y_0,y_1),
\]
where \(x_0,x_1\) belong to the first factor and \(y_0,y_1\) to the
second. A maximal cone contains one ray from each opposite pair. The
four complementary monomials are
\[
x_0y_0,
\qquad
x_0y_1,
\qquad
x_1y_0,
\qquad
x_1y_1.
\]
Hence
\[
\begin{aligned}
B(\Sigma)
&=
\langle x_0y_0,x_0y_1,x_1y_0,x_1y_1\rangle
\\
&=
\langle x_0,x_1\rangle
\cap
\langle y_0,y_1\rangle,
\end{aligned}
\]
and therefore
\[
\boxed{
Z(\Sigma)
=
V(x_0,x_1)
\cup
V(y_0,y_1).
}
\]

This is exactly the expected coordinate rule:
\[
(x_0,x_1)\neq(0,0)
\qquad\text{and}\qquad
(y_0,y_1)\neq(0,0).
\]
The two primitive collections are the two opposite pairs of rays.

\begin{pauseandthink}
Why is the origin not the only point that must be removed from
\(\AAA^4\) in this example? Give a point other than the origin that
cannot represent an element of \(\PP^1\times\PP^1\), and identify the
primitive collection responsible for excluding it.
\end{pauseandthink}

\subsection{Surface laboratory: the blowup of \texorpdfstring{\(\PP^2\)}{P2}}

For
\[
X=\Blp{\PP^2}{p},
\]
recall the cyclically ordered rays
\[
u_0=-e_1-e_2,
\qquad
u_1=e_1,
\qquad
u_E=e_1+e_2,
\qquad
u_2=e_2,
\]
with Cox variables
\[
x_0,x_1,x_E,x_2.
\]
The maximal cones and their complementary monomials are
\[
\begin{array}{c|c}
\text{maximal cone}&x^{\widehat\sigma}\\
\hline
\Cone{u_0,u_1}&x_Ex_2\\
\Cone{u_1,u_E}&x_0x_2\\
\Cone{u_E,u_2}&x_0x_1\\
\Cone{u_2,u_0}&x_1x_E
\end{array}
\]
Consequently,
\[
\begin{aligned}
B(\Sigma)
&=
\langle x_Ex_2,x_0x_2,x_0x_1,x_1x_E\rangle
\\
&=
\langle x_0,x_E\rangle
\cap
\langle x_1,x_2\rangle,
\end{aligned}
\]
so
\[
\boxed{
Z(\Sigma)
=
V(x_0,x_E)
\cup
V(x_1,x_2).
}
\]
The primitive collections are
\[
\{\rho_0,\rho_E\}
\qquad\text{and}\qquad
\{\rho_1,\rho_2\}.
\]
These are exactly the two pairs of rays that do not lie together in a
cone.

\begin{figure}[ht]
\centering
\begin{tikzpicture}[scale=0.95]
  % Fan panel
  \begin{scope}[xshift=-4.0cm]
    \node[font=\bfseries] at (0,3.0) {the fan};

    \fill[DictionaryBlue!8] (0,0)--(2.1,0)--(1.6,1.6)--cycle;
    \fill[PredictionTeal!8] (0,0)--(1.6,1.6)--(0,2.1)--cycle;
    \fill[ExampleGold!10] (0,0)--(0,2.1)--(-1.65,-1.65)--cycle;
    \fill[SurfaceGreen!9] (0,0)--(-1.65,-1.65)--(2.1,0)--cycle;

    \draw[->,very thick,DictionaryBlue] (0,0)--(2.3,0)
      node[right] {\(u_1\;[x_1]\)};
    \draw[->,very thick,PredictionTeal] (0,0)--(1.85,1.85)
      node[above right] {\(u_E\;[x_E]\)};
    \draw[->,very thick,ExampleGold] (0,0)--(0,2.35)
      node[above] {\(u_2\;[x_2]\)};
    \draw[->,very thick,SurfaceGreen] (0,0)--(-1.85,-1.85)
      node[below left] {\(u_0\;[x_0]\)};
    \fill (0,0) circle (1.8pt);

    \draw[WarningRed,dashed,thick,bend left=15]
      (-1.55,-1.55) to (1.45,1.45);
    \draw[WarningRed,dashed,thick,bend right=20]
      (2.0,0) to (0,2.0);
    \node[WarningRed,align=center,font=\small] at (-0.25,-2.45)
      {dashed pairs do not\\belong to one cone};
  \end{scope}

  % Exceptional set panel
  \begin{scope}[xshift=4.1cm]
    \node[font=\bfseries] at (0,3.0) {the forbidden coordinate planes};

    \draw[rounded corners,fill=WarningRed!7,draw=WarningRed,very thick]
      (-2.75,0.62) rectangle (2.75,2.05);
    \node at (0,1.57) {\(V(x_0,x_E)\)};
    \node[font=\scriptsize,align=center,text width=4.6cm] at (0,1.04)
      {the non-cone pair \(\{\rho_0,\rho_E\}\) vanishes};

    \node[WarningRed,font=\Large] at (0,0.22) {\(\cup\)};

    \draw[rounded corners,fill=WarningRed!7,draw=WarningRed,very thick]
      (-2.75,-1.57) rectangle (2.75,-0.14);
    \node at (0,-0.62) {\(V(x_1,x_2)\)};
    \node[font=\scriptsize,align=center,text width=4.6cm] at (0,-1.15)
      {the non-cone pair \(\{\rho_1,\rho_2\}\) vanishes};

    \node[align=center,font=\small] at (0,-2.28)
      {outside these planes, the zero coordinates\\fit inside at least one cone};
  \end{scope}
\end{tikzpicture}
\caption{For \(\Blp{\PP^2}{p}\), the fan has two primitive
collections. Each gives one forbidden coordinate plane in Cox space.}
\label{fig:blowup-irrelevant-exceptional-set}
\end{figure}

Notice how the blowup changes more than the grading. Inserting the ray
\(\rho_E\) changes which rays are adjacent, hence changes which
coordinate pairs may vanish simultaneously. The fan is now visible in
the open subset
\[
\AAA^4\setminus Z(\Sigma).
\]

\begin{prediction}
The open set \(D(x^{\widehat\sigma})\) has coordinate ring
\[
(S_\Sigma)_{x^{\widehat\sigma}},
\]
where every variable outside \(\sigma\) is invertible. Which monomials
in this localization should be unchanged by all class-group
rescalings? Predict why the degree-zero subring should be
\[
\Bbbk[\sigma^\vee\cap M],
\]
the coordinate ring of the affine toric chart \(U_\sigma\).
\end{prediction}

\begin{guidingquestion}
We now have two ingredients:
\[
\AAA^{\Sigma(1)}\setminus Z(\Sigma)
\qquad\text{and}\qquad
G=\Hom(\Cl(X_\Sigma),\Gm).
\]
Why does quotienting each open set \(D(x^{\widehat\sigma})\) by the
class-group rescalings recover \(U_\sigma\), and when are the fibers of
the resulting map exactly the \(G\)-orbits?
\end{guidingquestion}

\section{The quotient construction}
\label{sec:cox-quotient-construction}
\index{Cox quotient}
\index{rescaling group}
\index{degree-zero localization}
\index{quotient construction!Cox}

The previous two sections produced the ingredients of a homogeneous
coordinate construction:
\[
S_\Sigma=\Bbbk[x_\rho:\rho\in\Sigma(1)],
\qquad
\deg(x_\rho)=[D_\rho],
\qquad
\widehat X_\Sigma
=
\AAA^{\Sigma(1)}\setminus Z(\Sigma).
\]
The open set \(\widehat X_\Sigma\) remembers which coordinate patterns
are allowed by the fan. What remains is to identify points that differ
only by the rescalings encoded by the class-group grading.

Throughout this section, the rays of \(\Sigma\) span \(N_\RR\), so
\(X_\Sigma\) has no torus factor.

\subsection{The rescaling group}

The grading group itself determines the group that acts on Cox space.

\begin{definition}[The Cox rescaling group]
\label{def:cox-rescaling-group}
Set
\[
G
=
\Hom_{\ZZ}(\Cl(X_\Sigma),\Gm).
\]
For \(g\in G\), the value \(g(\alpha)\in\Gm\) is defined for every
class \(\alpha\in\Cl(X_\Sigma)\). The action on Cox coordinates is
\[
g\cdot(z_\rho)_\rho
=
\bigl(g([D_\rho])z_\rho\bigr)_\rho.
\]
\end{definition}

Thus a coordinate is rescaled according to its degree. More generally,
if \(f\in(S_\Sigma)_\alpha\) is homogeneous, then
\[
f(g\cdot z)=g(\alpha)f(z).
\]
In particular,
\[
\boxed{
\text{the }G\text{-invariant elements of }S_\Sigma
\text{ are precisely the elements of degree }0.
}
\]

\begin{observation}
If
\[
\Cl(X_\Sigma)\simeq\ZZ^r\oplus A_{\mathrm{tors}},
\]
then \(G\) is a product of \((\Gm)^r\) with a finite diagonalizable
group. The technical name is a \emph{quasitorus}. When the class group
is free, as in all of our smooth surface examples, \(G\) is an
ordinary algebraic torus.
\end{observation}

Dualizing the class-group sequence
\[
0\longrightarrow M
\longrightarrow\ZZ^{\Sigma(1)}
\longrightarrow\Cl(X_\Sigma)
\longrightarrow0
\]
gives an exact sequence of diagonalizable groups
\[
\boxed{
1\longrightarrow G
\longrightarrow(\Gm)^{\Sigma(1)}
\longrightarrow T_N
\longrightarrow1.
}
\]
The middle group rescales every Cox coordinate independently. Its
subgroup \(G\) consists exactly of the rescalings that disappear when
we pass to the dense torus \(T_N\) of \(X_\Sigma\).

\noindent
\textit{Reference.}
This is Cox--Little--Schenck, exact sequence~(5.1.2) and
Lemma~5.1.1.

\begin{figure}[ht]
\centering
\begin{tikzpicture}[node distance=0.55cm,every node/.style={font=\small}]
  \node[draw=DictionaryBlue,very thick,rounded corners,
        fill=DictionaryBlue!8,minimum width=2.25cm,minimum height=1.0cm,
        align=center]
    (grading) {class-group grading\\\(\deg(x_\rho)=[D_\rho]\)};

  \node[draw=PredictionTeal,very thick,rounded corners,
        fill=PredictionTeal!8,minimum width=2.3cm,minimum height=1.0cm,
        align=center,right=of grading]
    (group) {rescaling group\\\(G=\Hom(\Cl(X),\Gm)\)};

  \node[draw=ExampleGold,very thick,rounded corners,
        fill=ExampleGold!12,minimum width=2.25cm,minimum height=1.0cm,
        align=center,right=of group]
    (action) {diagonal action\\\(z_\rho\mapsto g([D_\rho])z_\rho\)};

  \node[draw=SurfaceGreen,very thick,rounded corners,
        fill=SurfaceGreen!10,minimum width=2.15cm,minimum height=1.0cm,
        align=center,right=of action]
    (invariants) {invariant functions\\degree \(0\)};

  \draw[->,very thick] (grading)--(group);
  \draw[->,very thick] (group)--(action);
  \draw[->,very thick] (action)--(invariants);
\end{tikzpicture}
\caption{A grading and a diagonal group action are two descriptions of
the same symmetry. Quotient functions are the homogeneous expressions
whose total class-group degree is zero.}
\label{fig:grading-group-action-invariants}
\end{figure}

\begin{agcheatsheet}
Suppose an algebraic group \(G\) acts on an affine variety
\(Y=\Spec(A)\). The invariant ring
\[
A^G=\{f\in A:g\cdot f=f\text{ for every }g\in G\}
\]
defines the affine quotient
\[
Y\mathbin{//}G=\Spec(A^G).
\]
The double slash reminds us that this construction is governed by
invariant functions.

A \emph{geometric quotient} has the stronger property that its fibers
are exactly the \(G\)-orbits. A categorical or good quotient may put
several orbits in one fiber; in that case the invariant functions
cannot distinguish those orbits.

For Cox coordinates, the local quotient rings will therefore be the
degree-zero parts of suitable localizations of \(S_\Sigma\).
\end{agcheatsheet}

\subsection{Recovering one affine chart}

Fix a cone \(\sigma\in\Sigma\). Recall the monomial
\[
x^{\widehat\sigma}
=
\prod_{\rho\notin\sigma(1)}x_\rho.
\]
The corresponding principal open subset is
\[
\widehat U_\sigma
=
D(x^{\widehat\sigma})
\subseteq\widehat X_\Sigma.
\]
On this open set, every coordinate \(x_\rho\) with
\(\rho\notin\sigma(1)\) is invertible. Its coordinate ring is
\[
\Bbbk[\widehat U_\sigma]
=
(S_\Sigma)_{x^{\widehat\sigma}}.
\]

The prediction at the end of the previous section is now the local
heart of the quotient construction.

\begin{theorem}[The degree-zero localization is the affine toric chart]
\label{thm:degree-zero-localization-affine-chart}
For every \(\sigma\in\Sigma\), there is a natural isomorphism
\[
\boxed{
\left((S_\Sigma)_{x^{\widehat\sigma}}\right)_0
\simeq
\Bbbk[\sigma^\vee\cap M].
}
\]
Equivalently,
\[
\widehat U_\sigma\mathbin{//}G
\simeq
U_\sigma.
\]
\end{theorem}

\noindent
\textit{Reference.}
This is the local invariant-ring computation~(5.1.4) in the proof of
Cox--Little--Schenck, Theorem~5.1.11.

\begin{proof}
For \(m\in\sigma^\vee\cap M\), define
\[
\Phi_\sigma(\chi^m)
=
\prod_{\rho\in\Sigma(1)}
 x_\rho^{\langle m,u_\rho\rangle}.
\]
When \(\rho\in\sigma(1)\), the exponent is nonnegative because
\(m\in\sigma^\vee\). When \(\rho\notin\sigma(1)\), the exponent may
be negative, but \(x_\rho\) is invertible in
\((S_\Sigma)_{x^{\widehat\sigma}}\). Thus the expression is well
defined.

Its degree is zero, since its exponent vector is the principal-divisor
vector
\[
\bigl(\langle m,u_\rho\rangle\bigr)_\rho
\in\im(M\to\ZZ^{\Sigma(1)}).
\]
Hence \(\Phi_\sigma\) maps into the degree-zero subring.

Conversely, consider a degree-zero Laurent monomial
\[
x^a=\prod_\rho x_\rho^{a_\rho}
\in(S_\Sigma)_{x^{\widehat\sigma}}.
\]
We have \(a_\rho\geq0\) for \(\rho\in\sigma(1)\), since those
variables were not inverted. Degree zero means that the exponent
vector \(a\) lies in the kernel of
\[
\ZZ^{\Sigma(1)}\longrightarrow\Cl(X_\Sigma).
\]
By exactness, there is a unique \(m\in M\) such that
\[
a_\rho=\langle m,u_\rho\rangle
\qquad\text{for every }\rho.
\]
The inequalities on the coordinates belonging to \(\sigma\) say
precisely that \(m\in\sigma^\vee\). Therefore every degree-zero
monomial is in the image of \(\Phi_\sigma\), and the same follows for
linear combinations of such monomials.
\end{proof}

The formula is worth reading in both directions:
\[
\boxed{
\chi^m
\quad\longleftrightarrow\quad
\prod_\rho x_\rho^{\langle m,u_\rho\rangle}.
}
\]
A character on the toric chart becomes a degree-zero Laurent monomial
in Cox coordinates. Negative exponents are allowed only on the
coordinates that were declared nonzero on that chart.

\begin{example}[The usual affine chart of projective space]
For \(\PP^2\), let \(\sigma\) be the cone whose complementary ray
corresponds to \(x_0\). Then
\[
\widehat U_\sigma=D(x_0),
\qquad
(S_{x_0})_0
=
\Bbbk\left[\frac{x_1}{x_0},\frac{x_2}{x_0}\right].
\]
Thus
\[
D(x_0)\mathbin{//}\Gm
\simeq
\AAA^2,
\]
which is the familiar projective chart \(x_0\neq0\).
\end{example}

\subsection{Gluing the local quotients}

If \(\tau\preceq\sigma\), then
\[
\widehat U_\tau\subseteq\widehat U_\sigma.
\]
Algebraically, passing from \(\sigma\) to \(\tau\) inverts additional
Cox variables. Under
\cref{thm:degree-zero-localization-affine-chart}, this is exactly the
usual localization
\[
\Bbbk[\sigma^\vee\cap M]
\longrightarrow
\Bbbk[\tau^\vee\cap M]
\]
that defines the open immersion \(U_\tau\hookrightarrow U_\sigma\).
The local quotients therefore glue in precisely the same pattern as
the affine toric charts.

\begin{theorem}[The Cox quotient]
\label{thm:cox-quotient-construction}
Let \(X_\Sigma\) be a toric variety without torus factors. There is a
surjective toric morphism
\[
\pi:\widehat X_\Sigma
=
\AAA^{\Sigma(1)}\setminus Z(\Sigma)
\longrightarrow
X_\Sigma
\]
that is constant on \(G\)-orbits, and
\[
\boxed{
X_\Sigma
\simeq
\widehat X_\Sigma\mathbin{//}G.
}
\]
Moreover, the quotient is geometric---its fibers are exactly the
\(G\)-orbits---if and only if \(\Sigma\) is simplicial.
\end{theorem}

\noindent
\textit{Reference.}
The construction of \(\pi\) is Cox--Little--Schenck,
Proposition~5.1.9, and the quotient statement is their
Theorem~5.1.11.

\begin{proof}[Conceptual proof sketch]
On every \(G\)-stable open set \(\widehat U_\sigma\), invariant
functions are the degree-zero localized functions. By
\cref{thm:degree-zero-localization-affine-chart}, their spectrum is
\(U_\sigma\). These affine quotients agree on overlaps, so they glue
to a global morphism
\[
\pi:\widehat X_\Sigma\to X_\Sigma.
\]

The remaining issue is orbit structure. If \(\sigma\) is simplicial,
its ray generators are linearly independent over \(\RR\); this
prevents a nontrivial one-parameter subgroup of \(G\) from moving an
orbit toward a different orbit inside the same local fiber. Hence each
fiber is one orbit. For a nonsimplicial cone, a linear dependence among
its rays produces distinct orbits that invariant functions cannot
separate.
\end{proof}

\begin{warningbox}
The notation matters:
\[
X_\Sigma=\widehat X_\Sigma\mathbin{//}G
\]
is valid in general. Writing a single slash
\[
X_\Sigma=\widehat X_\Sigma/G
\]
as a literal orbit space is justified when \(\Sigma\) is simplicial.
For a smooth fan the action is even free, so the quotient behaves like
the most familiar projective-space construction.
\end{warningbox}

\subsection{Surface laboratory: one chart of the blowup}

For
\[
X=\Blp{\PP^2}{p},
\]
use the basis
\[
L=H-E,
\qquad
E
\]
of \(\Cl(X)\). Then
\[
\deg(x_0)=L+E,
\qquad
\deg(x_1)=\deg(x_2)=L,
\qquad
\deg(x_E)=E.
\]
Consequently,
\[
G\simeq(\Gm)^2
\]
acts by
\[
(\lambda,\mu)\cdot(x_0,x_1,x_E,x_2)
=
(\lambda\mu x_0,\lambda x_1,\mu x_E,\lambda x_2).
\]

Consider the cone
\[
\sigma=\Cone{u_1,u_E}.
\]
The complementary variables are \(x_0,x_2\), so
\[
\widehat U_\sigma=D(x_0x_2).
\]
Two degree-zero functions are
\[
s=\frac{x_1}{x_2},
\qquad
 t=\frac{x_Ex_2}{x_0}.
\]
They are invariant under both \(\lambda\) and \(\mu\), and the local
theorem gives
\[
\left((S_X)_{x_0x_2}\right)_0
=
\Bbbk[s,t].
\]
Hence
\[
\widehat U_\sigma\mathbin{//}(\Gm)^2
\simeq
\AAA^2.
\]

\begin{figure}[ht]
\centering
\begin{tikzpicture}[scale=0.92]
  % Fan panel
  \begin{scope}[xshift=-4.4cm]
    \node[font=\bfseries] at (0,2.8) {choose a cone};

    \fill[PredictionTeal!13]
      (0,0)--(2.0,0)--(1.55,1.55)--cycle;
    \draw[->,very thick,DictionaryBlue] (0,0)--(2.25,0)
      node[right] {\(u_1\)};
    \draw[->,very thick,PredictionTeal] (0,0)--(1.8,1.8)
      node[above right] {\(u_E\)};
    \draw[->,very thick,ExampleGold] (0,0)--(0,2.25)
      node[above] {\(u_2\)};
    \draw[->,very thick,SurfaceGreen] (0,0)--(-1.75,-1.75)
      node[below left] {\(u_0\)};
    \fill (0,0) circle (1.8pt);

    \node[align=center,font=\small] at (0,-2.35)
      {outside \(\sigma\): invert \(x_0,x_2\)};
  \end{scope}

  % Algebra panel
  \begin{scope}[xshift=3.7cm]
    \node[font=\bfseries] at (0.3,2.8) {take degree-zero functions};

    \node[draw=DictionaryBlue,very thick,rounded corners,
          fill=DictionaryBlue!8,minimum width=3.4cm,minimum height=0.9cm]
      (loc) at (0.3,1.55) {\(D(x_0x_2)\)};

    \node[draw=PredictionTeal,very thick,rounded corners,
          fill=PredictionTeal!8,minimum width=3.4cm,minimum height=1.05cm,
          align=center]
      (inv) at (0.3,0.05)
      {\(s=x_1/x_2\), \(t=x_Ex_2/x_0\)\\both have degree \(0\)};

    \node[draw=SurfaceGreen,very thick,rounded corners,
          fill=SurfaceGreen!10,minimum width=3.4cm,minimum height=0.9cm]
      (aff) at (0.3,-1.55) {\(\Spec\Bbbk[s,t]=\AAA^2\)};

    \draw[->,very thick] (loc)--(inv)
      node[midway,right,font=\scriptsize] {quotient by \((\Gm)^2\)};
    \draw[->,very thick] (inv)--(aff);
  \end{scope}
\end{tikzpicture}
\caption{One Cox chart of the blowup. The variables outside the chosen
cone are inverted, and the degree-zero Laurent monomials become the
ordinary coordinates on the affine toric chart.}
\label{fig:blowup-local-cox-quotient}
\end{figure}

The same calculation on the other three maximal cones gives four
copies of \(\AAA^2\). Their degree-zero localizations glue according to
the fan, recovering the blowup.

\begin{pauseandthink}
For \(\PP^1\times\PP^1\), use Cox coordinates
\[
(x_0,x_1,y_0,y_1)
\]
with bidegrees
\[
\deg(x_i)=(1,0),
\qquad
\deg(y_j)=(0,1).
\]
On the open set \(D(x_0y_0)\), find two degree-zero generators and
show that the local quotient is \(\AAA^2\). How do your generators
compare with the usual affine coordinates on the two factors?
\end{pauseandthink}

\begin{guidingquestion}
The quotient construction is now formal. What does it look like in
specific examples when we write down, side by side,
\[
\text{the grading},
\qquad
Z(\Sigma),
\qquad
G,
\qquad
\text{and the resulting quotient}?
\]
How can two polynomial rings with the same number of variables produce
different toric varieties because their gradings or exceptional sets
differ?
\end{guidingquestion}

\section{Four quotient constructions}
\label{sec:four-cox-quotient-examples}
\index{weight matrix}
\index{projective space!Cox quotient}
\index{Hirzebruch surface!Cox quotient}
\index{blowup!Cox quotient}

The quotient theorem becomes most useful when its ingredients are
written next to one another:
\[
\text{Cox ring},
\qquad
\text{grading},
\qquad
\text{exceptional set},
\qquad
\text{group action}.
\]
The examples in this section all begin with a polynomial ring, but
small changes in the grading and the deleted coordinate locus produce
different varieties.

\subsection{Weight matrices}

Suppose for the moment that
\[
\Cl(X_\Sigma)\simeq\ZZ^k
\]
and choose a basis. If the rays are
\(\rho_1,\dots,\rho_s\), write the degree of each Cox variable as a
column vector. The resulting matrix
\[
Q=
\begin{pmatrix}
|&&|\\
\deg(x_{\rho_1})&\cdots&\deg(x_{\rho_s})\\
|&&|
\end{pmatrix}
\in\operatorname{Mat}_{k\times s}(\ZZ)
\]
is called a \emph{weight matrix} for the chosen basis.

If
\[
g=(\lambda_1,\dots,\lambda_k)\in(\Gm)^k,
\]
and the \(j\)-th column of \(Q\) is
\((q_{1j},\dots,q_{kj})^\mathsf{T}\), then
\[
g\cdot x_{\rho_j}
=
\lambda_1^{q_{1j}}\cdots\lambda_k^{q_{kj}}x_{\rho_j}.
\]
Thus the weight matrix turns the grading immediately into the diagonal
action.

\begin{warningbox}
A weight matrix depends on the chosen basis of \(\Cl(X_\Sigma)\).
Changing the basis changes the displayed matrix and the parameters
used for \(G\), but not the quotient variety.
\end{warningbox}

\subsection{Projective space}

For \(\PP^n\), the Cox ring is
\[
S=\Bbbk[x_0,\dots,x_n],
\qquad
\Cl(\PP^n)=\ZZ H,
\]
and every variable has degree \(H\). Thus
\[
Q_{\PP^n}
=
\begin{pmatrix}
1&1&\cdots&1
\end{pmatrix}.
\]
The irrelevant ideal is
\[
B(\Sigma)=\langle x_0,\dots,x_n\rangle,
\]
so the exceptional set is only the origin. The group
\[
G=\Gm
\]
acts by simultaneous scaling:
\[
\lambda\cdot(x_0,\dots,x_n)
=
(\lambda x_0,\dots,\lambda x_n).
\]
Therefore the Cox quotient is exactly the classical construction
\[
\boxed{
\PP^n
=
\bigl(\AAA^{n+1}\setminus\{0\}\bigr)/\Gm.
}
\]

\noindent
\textit{Reference.}
This is Cox--Little--Schenck, Example~5.1.12.

On the chart \(x_i\neq0\), the invariant functions are
\[
\frac{x_0}{x_i},\dots,
\widehat{\frac{x_i}{x_i}},\dots,
\frac{x_n}{x_i},
\]
recovering the usual affine chart \(\AAA^n\).

\subsection{A product: \texorpdfstring{\(\PP^1\times\PP^1\)}{P1 x P1}}

Write the Cox coordinates as
\[
(u_0,u_1,v_0,v_1).
\]
The class group and grading are
\[
\Cl(\PP^1\times\PP^1)
\simeq
\ZZ F_1\oplus\ZZ F_2,
\]
\[
\deg(u_0)=\deg(u_1)=F_1,
\qquad
\deg(v_0)=\deg(v_1)=F_2.
\]
Hence
\[
Q_{\PP^1\times\PP^1}
=
\begin{pmatrix}
1&1&0&0\\
0&0&1&1
\end{pmatrix}.
\]

The opposite pairs of rays are the primitive collections, so
\[
B(\Sigma)
=
\langle u_0,u_1\rangle
\cap
\langle v_0,v_1\rangle
\]
and
\[
Z(\Sigma)
=
V(u_0,u_1)
\cup
V(v_0,v_1).
\]
Thus an allowed point satisfies
\[
(u_0,u_1)\neq(0,0),
\qquad
(v_0,v_1)\neq(0,0).
\]
The group \((\Gm)^2\) acts by
\[
(\lambda,\mu)\cdot(u_0,u_1,v_0,v_1)
=
(\lambda u_0,\lambda u_1,
 \mu v_0,\mu v_1).
\]
The two blocks are independently projectivized, giving
\[
\boxed{
\PP^1\times\PP^1
=
\frac{
\AAA^4\setminus
\bigl(V(u_0,u_1)\cup V(v_0,v_1)\bigr)
}{(\Gm)^2}.
}
\]

\noindent
\textit{Reference.}
This is Cox--Little--Schenck, Example~5.1.13.

\begin{pauseandthink}
On the open set \(D(u_0v_0)\), check that
\[
s=\frac{u_1}{u_0},
\qquad
t=\frac{v_1}{v_0}
\]
have degree zero. Explain why they give the usual affine coordinates
on one product chart.
\end{pauseandthink}

\subsection{A twist: the Hirzebruch surface \texorpdfstring{\(\mathbb F_r\)}{F-r}}

Recall the fan of \(\mathbb F_r\), with rays in cyclic order
\[
w_1=(1,0),
\qquad
w_2=(0,1),
\qquad
w_3=(-1,r),
\qquad
w_4=(0,-1).
\]
Use Cox coordinates
\[
(u_0,v_0,u_1,v_1)
\]
in the same cyclic order. The principal-divisor relations are
\[
D_1-D_3\sim0,
\qquad
D_2+rD_3-D_4\sim0.
\]
Set
\[
F=[D_1]=[D_3],
\qquad
S=[D_2].
\]
Then
\[
[D_4]=S+rF,
\]
so the Cox grading is
\[
\deg(u_0)=\deg(u_1)=F,
\qquad
\deg(v_0)=S,
\qquad
\deg(v_1)=S+rF.
\]
With respect to the ordered basis \((F,S)\),
\[
Q_{\mathbb F_r}
=
\begin{pmatrix}
1&0&1&r\\
0&1&0&1
\end{pmatrix}.
\]
Consequently,
\[
(\lambda,\mu)\cdot(u_0,v_0,u_1,v_1)
=
(\lambda u_0,
 \mu v_0,
 \lambda u_1,
 \lambda^r\mu v_1).
\]

The fan is still a four-cycle, so its primitive collections are the
same opposite pairs as for \(\PP^1\times\PP^1\). Hence
\[
Z(\Sigma_r)
=
V(u_0,u_1)
\cup
V(v_0,v_1).
\]
We therefore obtain
\[
\boxed{
\mathbb F_r
=
\frac{
\AAA^4\setminus
\bigl(V(u_0,u_1)\cup V(v_0,v_1)\bigr)
}{(\Gm)^2},
}
\]
but with the twisted action displayed above.

\begin{figure}[ht]
\centering
\begin{tikzpicture}[scale=0.94]
  % Product panel
  \begin{scope}[xshift=-4.35cm]
    \node[font=\bfseries] at (0,2.65)
      {\(\PP^1\times\PP^1\)};

    \node[draw=DictionaryBlue,very thick,rounded corners,
      fill=DictionaryBlue!8,minimum width=5.3cm,minimum height=0.8cm]
      at (0,1.62)
      {delete \(V(u_0,u_1)\cup V(v_0,v_1)\)};

    \node[draw=PredictionTeal,very thick,rounded corners,
      fill=PredictionTeal!8,minimum width=5.3cm,minimum height=1.05cm,
      align=center]
      at (0,0.25)
      {\((\lambda,\mu)\):\\
       \((u_0,u_1)\mapsto\lambda(u_0,u_1)\),\\[-1mm]
       \((v_0,v_1)\mapsto\mu(v_0,v_1)\)};

    \node[align=center,font=\small] at (0,-1.55)
      {two independent\\untwisted scalings};
  \end{scope}

  % Hirzebruch panel
  \begin{scope}[xshift=4.35cm]
    \node[font=\bfseries] at (0,2.65)
      {\(\mathbb F_r\)};

    \node[draw=DictionaryBlue,very thick,rounded corners,
      fill=DictionaryBlue!8,minimum width=5.3cm,minimum height=0.8cm]
      at (0,1.62)
      {delete \(V(u_0,u_1)\cup V(v_0,v_1)\)};

    \node[draw=ExampleGold,very thick,rounded corners,
      fill=ExampleGold!10,minimum width=5.3cm,minimum height=1.05cm,
      align=center]
      at (0,0.25)
      {\((\lambda,\mu)\):\\
       \((u_0,u_1)\mapsto\lambda(u_0,u_1)\),\\[-1mm]
       \((v_0,v_1)\mapsto(\mu v_0,\lambda^r\mu v_1)\)};

    \node[align=center,font=\small] at (0,-1.55)
      {the same deleted set,\\but one coordinate twists by \(\lambda^r\)};
  \end{scope}
\end{tikzpicture}
\caption{The product and the Hirzebruch surface use the same ambient
space and the same exceptional set. The difference lies entirely in
the grading, hence in the diagonal action.}
\label{fig:product-versus-hirzebruch-cox-action}
\end{figure}

The quotient map to the base is visible immediately:
\[
\pi:\mathbb F_r\longrightarrow\PP^1,
\qquad
[u_0,v_0,u_1,v_1]
\longmapsto
[u_0:u_1].
\]
On the chart \(u_0\neq0\), use \(\lambda\) to set \(u_0=1\). The
remaining pair of fiber coordinates is
\[
[v_0:u_0^{-r}v_1].
\]
On the chart \(u_1\neq0\), it is
\[
[v_0:u_1^{-r}v_1].
\]
On the overlap, these differ by the factor
\[
\left(\frac{u_1}{u_0}\right)^{-r}.
\]
The exponent \(r\) in the Cox action is therefore exactly the twisting
seen when the two local product charts are glued.

When \(r=0\), the twist disappears and
\[
\mathbb F_0\simeq\PP^1\times\PP^1.
\]

\begin{prediction}
The exceptional set does not depend on \(r\), but the quotient variety
does. Which part of the quotient data must therefore remember the
self-intersection of the distinguished section of \(\mathbb F_r\)?
\end{prediction}

\subsection{The blowup of \texorpdfstring{\(\PP^2\)}{P2}}

For
\[
X=\Blp{\PP^2}{p},
\]
use the basis
\[
L=H-E,
\qquad E
\]
of \(\Cl(X)\). The degrees from
\cref{sec:why-homogeneous-coordinates-need-grading} become
\[
\deg(x_0)=L+E,
\qquad
\deg(x_1)=\deg(x_2)=L,
\qquad
\deg(x_E)=E.
\]
In the coordinate order
\[
(x_0,x_1,x_E,x_2),
\]
the weight matrix is
\[
Q_X
=
\begin{pmatrix}
1&1&0&1\\
1&0&1&0
\end{pmatrix}.
\]
Thus
\[
(\lambda,\mu)\cdot(x_0,x_1,x_E,x_2)
=
(\lambda\mu x_0,
 \lambda x_1,
 \mu x_E,
 \lambda x_2).
\]
From \cref{sec:irrelevant-ideal-exceptional-set},
\[
Z(\Sigma)
=
V(x_0,x_E)
\cup
V(x_1,x_2).
\]
Hence
\[
\boxed{
\Blp{\PP^2}{p}
=
\frac{
\AAA^4\setminus
\bigl(V(x_0,x_E)\cup V(x_1,x_2)\bigr)
}{(\Gm)^2}.
}
\]

This is the \(r=1\) Hirzebruch quotient in different names. Indeed,
under the relabeling
\[
\begin{array}{c|cccc}
\mathbb F_1&u_0&u_1&v_0&v_1\\
\hline
\Blp{\PP^2}{p}&x_1&x_2&x_E&x_0,
\end{array}
\]
and the identification
\[
F=L,
\qquad
S=E,
\]
both the degree matrix and the exceptional set agree. Thus Cox
coordinates make the isomorphism
\[
\boxed{
\mathbb F_1\simeq\Blp{\PP^2}{p}
}
\]
visible as a simple relabeling of variables and divisor classes.

\begin{surfacelaboratory}
The four examples can be compared at a glance:

\medskip

\begin{center}
\small
\begin{tabularx}{\linewidth}{@{}
  >{\raggedright\arraybackslash}p{0.16\linewidth}
  >{\raggedright\arraybackslash}p{0.22\linewidth}
  >{\raggedright\arraybackslash}X
  >{\raggedright\arraybackslash}p{0.21\linewidth}@{}}
\toprule
\textbf{Variety}
&
\textbf{Cox group}
&
\textbf{Deleted locus}
&
\textbf{Decisive feature}
\\
\midrule
\(\PP^n\)
&
\(\Gm\)
&
origin
&
one common scaling
\\
\(\PP^1\times\PP^1\)
&
\((\Gm)^2\)
&
\(V(u_0,u_1)\cup V(v_0,v_1)\)
&
two untwisted blocks
\\
\(\mathbb F_r\)
&
\((\Gm)^2\)
&
the same two planes
&
weight \(r\) twists one coordinate
\\
\(\Blp{\PP^2}{p}\)
&
\((\Gm)^2\)
&
\(V(x_0,x_E)\cup V(x_1,x_2)\)
&
new ray gives a new coordinate and class
\\
\bottomrule
\end{tabularx}
\end{center}
\end{surfacelaboratory}

The examples support a useful slogan:
\[
\boxed{
\begin{gathered}
\text{The fan determines what must be deleted}\\
\text{and how the remaining coordinates must be rescaled.}
\end{gathered}
}
\]
The irrelevant ideal remembers which rays fit into cones, while the
class-group grading remembers the linear relations among those rays.
Neither piece alone is enough.

\begin{guidingquestion}
All four fans in this section are smooth, so the Cox group acts freely
on the Cox open set and the quotient is geometric. What changes when a
cone is simplicial but not smooth, or when it is not simplicial at all?

Should one expect finite stabilizers, positive-dimensional
stabilizers, or fibers containing more than one orbit?
\end{guidingquestion}
\section{Singular and non-simplicial fans}
\label{sec:singular-nonsimplicial-cox-quotients}
\index{stabilizer!Cox action}
\index{simplicial fan}
\index{nonsimplicial fan}
\index{cyclic quotient singularity}

The examples in \cref{sec:four-cox-quotient-examples} were smooth.
Their Cox groups acted freely on the Cox open set, so passing to the
quotient simply identified points related by rescaling.

For a singular toric variety, two new phenomena can occur:
\[
\text{a point may have a nontrivial stabilizer},
\]
and
\[
\text{one quotient fiber may contain several distinct orbits}.
\]
The fan determines exactly which phenomenon occurs. The decisive
question is whether a cone is smooth, merely simplicial, or not
simplicial at all.

\subsection{Which rescalings fix a Cox point?}

Let
\[
z=(z_\rho)_\rho\in\widehat X_\Sigma
\]
be a point of the Cox open set, and let
\[
J(z)=\{\rho\in\Sigma(1):z_\rho\neq0\}.
\]
The coordinates indexed by \(J(z)\) are the ones that constrain a
rescaling: if \(z_\rho\neq0\), then a group element fixing \(z\) must
act trivially on that coordinate.

Set
\[
L(z)=
\left\langle [D_\rho]:\rho\in J(z)\right\rangle
\subseteq\Cl(X_\Sigma).
\]

\begin{proposition}[Stabilizers from surviving degrees]
\label{prop:cox-stabilizer-degree-quotient}
The stabilizer of \(z\) under the Cox group
\[
G=\Hom_{\ZZ}(\Cl(X_\Sigma),\Gm)
\]
is
\[
\boxed{
G_z
\simeq
\Hom_{\ZZ}\bigl(\Cl(X_\Sigma)/L(z),\Gm\bigr).
}
\]
\end{proposition}

\begin{proof}
An element \(g\in G\) fixes \(z\) precisely when
\[
g([D_\rho])z_\rho=z_\rho
\]
for every ray \(\rho\). When \(z_\rho=0\), this imposes no condition.
When \(z_\rho\neq0\), it is equivalent to
\[
g([D_\rho])=1.
\]
Thus \(g\) fixes \(z\) exactly when it is trivial on the subgroup
\(L(z)\). Such homomorphisms are precisely those that factor through
\(\Cl(X_\Sigma)/L(z)\).
\end{proof}

The formula becomes especially geometric at one distinguished point
for each cone. For \(\sigma\in\Sigma\), define
\[
z_\sigma=(z_\rho)_\rho,
\qquad
z_\rho=
\begin{cases}
0,&\rho\in\sigma(1),\\
1,&\rho\notin\sigma(1).
\end{cases}
\]
The zero coordinates fit inside \(\sigma\), so
\(z_\sigma\in\widehat X_\Sigma\). Its image in \(X_\Sigma\) lies in
the orbit \(O(\sigma)\).

The boundary divisors \(D_\rho\) with \(\rho\notin\sigma(1)\) do not
meet the affine chart \(U_\sigma\). Restriction therefore gives
\[
\Cl(U_\sigma)
\simeq
\frac{\Cl(X_\Sigma)}
{\left\langle[D_\rho]:\rho\notin\sigma(1)\right\rangle}.
\]
This is the toric instance of the divisor-class restriction sequence
of Cox--Little--Schenck, Theorem~4.0.20. Equivalently, one can apply
\cref{thm:toric-class-group-sequence} directly to the affine fan of
\(U_\sigma\).

\begin{corollary}[The local class group is the stabilizer group]
\label{cor:local-class-group-cox-stabilizer}
For every cone \(\sigma\in\Sigma\),
\[
\boxed{
G_{z_\sigma}
\simeq
\Hom_{\ZZ}\bigl(\Cl(U_\sigma),\Gm\bigr).
}
\]
\end{corollary}

This gives a useful interpretation of the local class group:
\[
\boxed{
\Cl(U_\sigma)
\text{ measures the rescaling symmetry left over above }O(\sigma).
}
\]

\subsection{The smooth--simplicial--nonsimplicial trichotomy}

Let \(\sigma(1)=\{\rho_1,\dots,\rho_s\}\), with primitive ray
generators \(u_1,\dots,u_s\). The affine class-group sequence is
\[
M
\longrightarrow
\ZZ^s
\longrightarrow
\Cl(U_\sigma)
\longrightarrow0,
\qquad
m\longmapsto
\bigl(\langle m,u_1\rangle,\dots,
      \langle m,u_s\rangle\bigr).
\]
The image has rank \(\dim\sigma\). Hence
\[
\rank\Cl(U_\sigma)
=
|\sigma(1)|-\dim\sigma.
\]

\begin{theorem}[Three quotient regimes]
\label{thm:three-cox-stabilizer-regimes}
For a cone \(\sigma\in\Sigma\), the following three cases occur.

\begin{enumerate}[label=(\arabic*),wide=0pt]
\item If \(\sigma\) is smooth, then
\[
\Cl(U_\sigma)=0,
\qquad
G_{z_\sigma}=\{1\}.
\]
The Cox action is free near the corresponding orbit.

\item If \(\sigma\) is simplicial but not smooth, then
\[
\Cl(U_\sigma)
\]
is a nonzero finite group, and \(G_{z_\sigma}\) is finite. The quotient
is geometric, but it has finite quotient singularities.

\item If \(\sigma\) is not simplicial, then
\[
\rank\Cl(U_\sigma)>0,
\]
so \(G_{z_\sigma}\) contains a positive-dimensional torus. The Cox
quotient is not geometric: invariant functions may identify several
orbits in one fiber.
\end{enumerate}
\end{theorem}

\noindent
\textit{Reference.}
The geometric-quotient criterion is Cox--Little--Schenck,
Theorem~5.1.11. The equivalence between simplicial fans and finite
quotient singularities is their Theorem~3.1.19. The divisor-theoretic
counterpart is \cref{thm:factorial-qfactorial-fan}.

\begin{proof}
The rank formula shows that \(\Cl(U_\sigma)\) is finite exactly when
\[
|\sigma(1)|=\dim\sigma,
\]
which is precisely the simplicial condition. If the cone is not
simplicial, the local class group has positive rank, so its dual
contains a positive-dimensional torus.

When \(\sigma\) is simplicial, the group is trivial precisely when the
map
\[
M\longrightarrow\ZZ^{\sigma(1)}
\]
is surjective. Dually, this says that the primitive ray generators
form part of a lattice basis, which is exactly the smoothness
condition. Now apply
\cref{cor:local-class-group-cox-stabilizer} and the quotient theorem
\cref{thm:cox-quotient-construction}.
\end{proof}

\begin{figure}[ht]
\centering
\begin{tikzpicture}[node distance=0.55cm,every node/.style={font=\small}]
  \node[draw=SurfaceGreen,very thick,rounded corners,
        fill=SurfaceGreen!9,text width=2.75cm,minimum height=0.9cm,
        align=center]
    (smooth) {smooth cone\\ray generators extend to a basis};
  \node[draw=SurfaceGreen,very thick,rounded corners,
        fill=SurfaceGreen!5,text width=2.75cm,minimum height=0.8cm,
        align=center,below=of smooth]
    (trivial) {trivial stabilizer\\free action};

  \node[draw=ExampleGold,very thick,rounded corners,
        fill=ExampleGold!12,text width=2.75cm,minimum height=0.9cm,
        align=center,right=0.65cm of smooth]
    (simplicial) {simplicial, not smooth\\independent over \(\RR\), not a basis};
  \node[draw=ExampleGold,very thick,rounded corners,
        fill=ExampleGold!7,text width=2.75cm,minimum height=0.8cm,
        align=center,below=of simplicial]
    (finite) {finite stabilizer\\geometric quotient};

  \node[draw=WarningRed,very thick,rounded corners,
        fill=WarningRed!7,text width=2.75cm,minimum height=0.9cm,
        align=center,right=0.65cm of simplicial]
    (nonsimplicial) {nonsimplicial cone\\more rays than its dimension};
  \node[draw=WarningRed,very thick,rounded corners,
        fill=WarningRed!4,text width=2.75cm,minimum height=0.8cm,
        align=center,below=of nonsimplicial]
    (positive) {positive-dimensional stabilizer\\several orbits may share a fiber};

  \draw[->,very thick,SurfaceGreen] (smooth)--(trivial);
  \draw[->,very thick,ExampleGold] (simplicial)--(finite);
  \draw[->,very thick,WarningRed] (nonsimplicial)--(positive);
\end{tikzpicture}
\caption{The local geometry of a toric chart is reflected in the
stabilizer of its distinguished Cox point.}
\label{fig:three-cox-stabilizer-regimes}
\end{figure}

\subsection{Finite stabilizers: a cyclic quotient surface}

Work over \(\Bbbk=\CC\), and let
\[
\sigma_d
=
\Cone{(d,-1),(0,1)}
\subseteq\RR^2,
\qquad d\geq2.
\]
The cone has two rays in dimension two, so it is simplicial. The
determinant of its primitive generators has absolute value \(d\), so
it is not smooth.

Let the Cox variables be \(x_1,x_2\). Since the fan has a single
maximal cone containing both rays,
\[
B(\Sigma)=\langle1\rangle,
\qquad
Z(\Sigma)=\varnothing.
\]
The class-group computation gives
\[
\Cl(U_{\sigma_d})\simeq\ZZ/d\ZZ,
\qquad
\deg(x_1)=\deg(x_2)=1\pmod d.
\]
Consequently,
\[
G=\Hom(\ZZ/d\ZZ,\Gm)=\mu_d,
\]
and a primitive \(d\)-th root of unity acts by
\[
\zeta\cdot(x_1,x_2)
=
(\zeta x_1,\zeta x_2).
\]
Thus
\[
\boxed{
U_{\sigma_d}
\simeq
\AAA^2/\mu_d.
}
\]

The invariant monomials are those whose total degree is divisible by
\(d\). In particular, the invariant ring is generated by
\[
x_1^d,
\quad
x_1^{d-1}x_2,
\quad\dots,\quad
x_1x_2^{d-1},
\quad
x_2^d.
\]
This is the affine cone over the degree-\(d\) rational normal curve.
For \(d=2\), setting
\[
A=x_1^2,
\qquad
B=x_1x_2,
\qquad
C=x_2^2
\]
gives
\[
\CC[x_1,x_2]^{\mu_2}
\simeq
\CC[A,B,C]/(AC-B^2).
\]

\noindent
\textit{Reference.}
The class-group calculation is Cox--Little--Schenck,
Example~4.1.4.

\begin{pauseandthink}
Away from the origin, determine which points of \(\AAA^2\) have a
nontrivial stabilizer under the scalar \(\mu_d\)-action. Why can the
quotient still be geometric even though the action is not free?
\end{pauseandthink}

The quotient map remembers each finite orbit as one point. The
singularity comes from the fact that the orbit of the origin has a
nontrivial finite stabilizer, not from several different orbits being
collapsed together.

\subsection{Positive-dimensional stabilizers: the cone over a square}

Now let \(N=\ZZ^3\), and consider the cone generated by
\[
\begin{aligned}
u_1&=(0,0,1),&
u_2&=(1,0,1),\\
u_3&=(1,1,1),&
u_4&=(0,1,1).
\end{aligned}
\]
These four rays are the cone from the origin over the vertices of a
unit square in the plane \(z=1\). The cone is three-dimensional but
has four rays, so it is not simplicial. The relation
\[
u_1-u_2+u_3-u_4=0
\]
shows that
\[
\Cl(U_\sigma)\simeq\ZZ,
\]
with Cox degrees
\[
\deg(x_1)=\deg(x_3)=1,
\qquad
\deg(x_2)=\deg(x_4)=-1.
\]
Hence \(G=\Gm\) acts by
\[
\lambda\cdot(x_1,x_2,x_3,x_4)
=
(\lambda x_1,\lambda^{-1}x_2,
 \lambda x_3,\lambda^{-1}x_4).
\]
Again the fan is affine with one maximal cone, so no locus is deleted:
\[
\widehat X_\Sigma=\AAA^4.
\]

The degree-zero invariant ring is generated by
\[
A=x_1x_2,
\qquad
B=x_1x_4,
\qquad
C=x_3x_2,
\qquad
D=x_3x_4.
\]
These generators satisfy one relation,
\[
AD=BC.
\]
Therefore
\[
\boxed{
U_\sigma
\simeq
\AAA^4\mathbin{//}\Gm
\simeq
\Spec\frac{\Bbbk[A,B,C,D]}{(AD-BC)}.
}
\]
This is the three-dimensional affine quadric cone.

\begin{figure}[ht]
\centering
\begin{tikzpicture}[scale=0.92]
  % cone over square
  \begin{scope}[xshift=-4.2cm]
    \coordinate (O) at (0,-1.35);
    \coordinate (P1) at (-1.45,0.3);
    \coordinate (P2) at (0.2,0.72);
    \coordinate (P3) at (1.5,0.12);
    \coordinate (P4) at (-0.15,-0.3);

    \fill[DictionaryBlue!9] (P1)--(P2)--(P3)--(P4)--cycle;
    \draw[DictionaryBlue,very thick] (P1)--(P2)--(P3)--(P4)--cycle;
    \foreach \P in {P1,P2,P3,P4} {
      \draw[ray] (O)--(\P);
      \fill (\P) circle (2pt);
    }
    \fill (O) circle (2pt);
    \node[below] at (O) {
      \(0\)};
    \node[above left] at (P1) {\(u_1\)};
    \node[above] at (P2) {\(u_2\)};
    \node[above right] at (P3) {\(u_3\)};
    \node[below left] at (P4) {\(u_4\)};
    \node[font=\bfseries] at (0,1.55) {four rays in dimension three};
    \node[align=center] at (0,-2.05)
      {nonsimplicial cone\\over a square};
  \end{scope}

  % quotient map and fiber
  \begin{scope}[xshift=4.2cm]
    \node[draw=PredictionTeal,very thick,rounded corners,
          fill=PredictionTeal!8,minimum width=3.1cm,minimum height=1.0cm,
          align=center] (a4) at (0,0.75)
      {\(\AAA^4\)\\weights \((1,-1,1,-1)\)};
    \node[draw=WarningRed,very thick,rounded corners,
          fill=WarningRed!6,minimum width=3.1cm,minimum height=1.0cm,
          align=center] (quadric) at (0,-1.05)
      {\(V(AD-BC)\subseteq\AAA^4\)\\three-dimensional quadric cone};
    \draw[->,very thick] (a4)--node[right] {\(//\Gm\)} (quadric);
  \end{scope}
\end{tikzpicture}
\caption{A non-simplicial cone produces a positive-dimensional
stabilizer. The Cox quotient of \(\AAA^4\) is the affine quadric cone,
and the fiber over its vertex is larger than one orbit.}
\label{fig:nonsimplicial-square-cone-quotient}
\end{figure}

The failure of a geometric quotient can be seen directly. The vertex
of the quadric cone is defined by
\[
A=B=C=D=0.
\]
Its inverse image is
\[
V(x_1,x_3)
\cup
V(x_2,x_4).
\]
For example, every point
\[
(x_1,0,x_3,0)\neq0
\]
maps to the vertex, but two such points need not lie in the same
\(\Gm\)-orbit: the ratio \([x_1:x_3]\in\PP^1\) is unchanged by the
action. Thus one quotient fiber contains a whole family of distinct
orbits, together with the origin whose stabilizer is all of \(\Gm\).

This is exactly what the double slash warns us about:
\[
\AAA^4\mathbin{//}\Gm
\]
is governed by invariant functions, and those functions cannot
separate the different orbits in the vertex fiber.

\begin{warningbox}
Two uses of the word ``quotient'' must be kept separate.

The \emph{Cox quotient}
\[
\widehat X_\Sigma\mathbin{//}G
\]
is a global presentation of every toric variety. The group \(G\) may
have positive dimension even when the toric variety is smooth; for
example, \(\PP^n\) is a quotient by \(\Gm\).

A \emph{finite quotient singularity} is a local statement: near a
singular point, the variety is locally a quotient of a smooth variety
by a finite group. For toric varieties, this happens precisely for
simplicial fans. What must be finite are the relevant point
stabilizers, not necessarily the entire Cox group.
\end{warningbox}

\begin{threedictionaries}
The singular quotient picture can now be read in three languages.

\medskip

\begin{center}
\small
\begin{tabularx}{\linewidth}{@{}
  >{\raggedright\arraybackslash}X
  >{\raggedright\arraybackslash}X
  >{\raggedright\arraybackslash}X@{}}
\toprule
\textbf{Algebra}
&
\textbf{Geometry}
&
\textbf{Combinatorics}
\\
\midrule
\(\Cl(U_\sigma)=0\)
&
trivial stabilizer; smooth chart
&
\(\sigma\) smooth
\\
\(\Cl(U_\sigma)\) finite
&
finite stabilizer; quotient singularity
&
\(\sigma\) simplicial
\\
\(\rank\Cl(U_\sigma)>0\)
&
positive-dimensional stabilizer; several orbits may share a fiber
&
\(\sigma\) nonsimplicial
\\
\bottomrule
\end{tabularx}
\end{center}
\end{threedictionaries}

\begin{guidingquestion}
The quotient construction began with a polynomial ring, a grading, and
an irrelevant ideal. How much of the original fan can be recovered
from those data alone?

More specifically, if one changes the deleted locus while keeping the
same graded polynomial ring and the same group action, what kinds of
different toric quotients might appear?
\end{guidingquestion}

\clearpage
\section*{What this chapter has accomplished}
\addcontentsline{toc}{section}{What this chapter has accomplished}
\markboth{What this chapter has accomplished}{What this chapter has accomplished}

A fan \(\Sigma\) with rays spanning \(N_\RR\) produces the polynomial
ring
\[
S_\Sigma=\Bbbk[x_\rho:\rho\in\Sigma(1)].
\]
The ring itself remembers only the rays. The class-group grading
\[
\deg(x_\rho)=[D_\rho]\in\Cl(X_\Sigma)
\]
records the relations among their primitive generators and turns
ordinary homogeneity into a multigraded notion.

For a torus-invariant divisor \(D\), the homogeneous piece of degree
\([D]\) recovers its global sections:
\[
\boxed{
(S_\Sigma)_{[D]}
\simeq
H^0(X_\Sigma,\mathcal O_{X_\Sigma}(D)).
}
\]
The homogenization formula
\[
\character{m}
\longmapsto
\prod_\rho
x_\rho^{\langle m,u_\rho\rangle+a_\rho},
\qquad
D=\sum_\rho a_\rho D_\rho,
\]
identifies a lattice point of \(P_D\) with a Cox monomial of degree
\([D]\).

The fan determines which coordinate patterns are allowed. Its
irrelevant ideal is
\[
B(\Sigma)
=
\left\langle
x^{\widehat\sigma}
=
\prod_{\rho\notin\sigma(1)}x_\rho:
\sigma\in\Sigma_{\max}
\right\rangle,
\]
and the exceptional set
\[
Z(\Sigma)=V(B(\Sigma))
\]
removes precisely the points whose zero coordinates do not fit inside
one cone.

The grading determines the rescaling group
\[
G=\Hom_{\ZZ}(\Cl(X_\Sigma),\Gm),
\]
and the toric variety is recovered as
\[
\boxed{
X_\Sigma
\simeq
\bigl(\AAA^{\Sigma(1)}\setminus Z(\Sigma)\bigr)
\mathbin{//}G.
}
\]
Locally, the construction is the degree-zero identity
\[
\left((S_\Sigma)_{x^{\widehat\sigma}}\right)_0
\simeq
\Bbbk[\sigma^\vee\cap M].
\]
Thus Cox coordinates do not replace the affine toric charts; they
place all of them inside one multigraded affine space.

Finally, the stabilizers of the Cox action detect singularities. At
the distinguished point associated to a cone \(\sigma\),
\[
G_{z_\sigma}
\simeq
\Hom_{\ZZ}(\Cl(U_\sigma),\Gm).
\]
Consequently,
\[
\begin{array}{c@{\qquad\Longleftrightarrow\qquad}c}
\sigma\text{ smooth}
&
G_{z_\sigma}\text{ trivial},
\\
\sigma\text{ simplicial}
&
G_{z_\sigma}\text{ finite},
\\
\sigma\text{ nonsimplicial}
&
G_{z_\sigma}\text{ positive-dimensional}.
\end{array}
\]
The challenge exercise uses the same quotient language to construct
line bundles by adding one more coordinate.

\begin{exercises}[Comprehension exercises]
\markboth{Comprehension exercises}{Comprehension exercises}

\item For \(\PP^n\), recover from the fan the sequence
\[
0\longrightarrow M
\longrightarrow\ZZ^{n+1}
\longrightarrow\ZZ
\longrightarrow0.
\]
Show that every Cox variable has degree \(1\), that
\[
B(\Sigma)=\langle x_0,\ldots,x_n\rangle,
\]
and that the Cox quotient is the usual homogeneous-coordinate
construction of projective space. Identify \(S_d\) with
\(H^0(\PP^n,\mathcal O(d))\).

\item For \(X=\PP^1\times\PP^1\), let
\[
S=\Bbbk[x_0,x_1,y_0,y_1]
\]
with the standard bidegree.
\begin{enumerate}[label=(\alph*),wide=0pt]
\item Determine all monomials of bidegree \((2,1)\).
\item Recover
\[
h^0(X,\mathcal O(2,1))=6.
\]
\item Compute the irrelevant ideal and exceptional set.
\item Verify on the chart \(x_0y_0\neq0\) that the degree-zero
localized ring is a polynomial ring in two variables.
\end{enumerate}

\item Let \(X=\Blp{\PP^2}{p}\), with Cox degrees
\[
\deg(x_0)=H,
\qquad
\deg(x_1)=\deg(x_2)=H-E,
\qquad
\deg(x_E)=E.
\]
\begin{enumerate}[label=(\alph*),wide=0pt]
\item Write the weight matrix for the action of \((\Gm)^2\).
\item Compute all monomials of degree \(2H-E\).
\item Recompute
\[
B(\Sigma)
=
\langle x_Ex_2,x_0x_2,x_0x_1,x_1x_E\rangle.
\]
\item On the open set \(D(x_0x_2)\), find two algebraically
independent degree-zero Laurent monomials and recover an affine-plane
chart.
\end{enumerate}

\item For the Hirzebruch surface \(\mathbb F_r\), use the quotient
presentation from \cref{sec:four-cox-quotient-examples}.
\begin{enumerate}[label=(\alph*),wide=0pt]
\item Check directly that the two pairs of forbidden simultaneous
vanishings are independent of \(r\).
\item Show that the exponent \(r\) appears in the transition function
between the two local product charts.
\item Explain how the grading distinguishes \(\mathbb F_r\) from
\(\PP^1\times\PP^1\), even though the ambient affine space and deleted
locus are the same.
\end{enumerate}

\item Let \(\Sigma\) be a fan and let
\[
C=\{\rho_1,\ldots,\rho_s\}\subseteq\Sigma(1).
\]
Prove directly from \cref{prop:cone-test-cox-coordinates} that
\[
V(x_{\rho_1},\ldots,x_{\rho_s})\subseteq Z(\Sigma)
\]
if and only if \(C\) is contained in no cone of \(\Sigma\). Deduce the
description of \(Z(\Sigma)\) as the union of the coordinate subspaces
associated to primitive collections.

\item Consider the simplicial cone
\[
\sigma_d=\Cone{de_1-e_2,e_2}\subseteq\RR^2.
\]
\begin{enumerate}[label=(\alph*),wide=0pt]
\item Compute \(\Cl(U_{\sigma_d})\).
\item Determine the corresponding finite Cox group.
\item Compute its action on \(\AAA^2\).
\item Explain why the quotient is geometric even though the origin
has a nontrivial stabilizer.
\end{enumerate}

\item For the action
\[
\lambda\cdot(x_1,x_2,x_3,x_4)
=
(\lambda x_1,\lambda^{-1}x_2,
 \lambda x_3,\lambda^{-1}x_4),
\]
show that the invariant ring is generated by
\[
A=x_1x_2,
\quad
B=x_1x_4,
\quad
C=x_3x_2,
\quad
D=x_3x_4
\]
with the single relation \(AD-BC=0\). Find two points in the fiber
over the vertex that are not in the same orbit, and explain why this
cannot happen for a geometric quotient.

\item Let \(z\in\widehat X_\Sigma\), and let
\[
I(z)=\{\rho:z_\rho\neq0\}.
\]
Show that the stabilizer of \(z\) is the subgroup of \(G\) on which
all characters \([D_\rho]\), \(\rho\in I(z)\), are trivial. Apply this
to a distinguished point \(z_\sigma\) and recover
\[
G_{z_\sigma}
\simeq
\Hom(\Cl(U_\sigma),\Gm).
\]

\item In \(N=\ZZ^3\), let
\[
u_1=(1,0,1),\quad
u_2=(0,1,1),\quad
u_3=(-1,0,1),\quad
u_4=(0,-1,1).
\]
The cone over the square with these vertices has two simplicial
subdivisions:
\[
\Sigma_+
=
\{\Cone{u_1,u_2,u_3},\Cone{u_1,u_3,u_4}\}
\]
and
\[
\Sigma_-
=
\{\Cone{u_1,u_2,u_4},\Cone{u_2,u_3,u_4}\},
\]
together with all faces. Compare their Cox constructions.
\begin{enumerate}[label=(\alph*),wide=0pt]
\item Show that the polynomial ring, class-group grading, and group
\(G\) are the same for both fans.
\item Compute the two irrelevant ideals and show that they differ.
\item Explain why the irrelevant ideal is essential for recovering the
whole fan rather than only its rays.
\end{enumerate}

\end{exercises}

\clearpage
\section*{Challenge exercise}
\addcontentsline{toc}{section}{Challenge exercise}
\markboth{Challenge exercise}{Challenge exercise}

\noindent
\textbf{Line bundles as an extra Cox coordinate.}

\medskip

Let \(X=X_\Sigma\) be a smooth toric variety without torus factors.
Write
\[
\widehat X
=
\AAA^{\Sigma(1)}\setminus Z(\Sigma),
\qquad
G=\Hom_{\ZZ}(\Cl(X),\Gm),
\]
and let
\[
\pi:\widehat X\longrightarrow X
\]
be the Cox quotient. Since \(X\) is smooth, the action of \(G\) on
\(\widehat X\) is free.

Every class \(\alpha\in\Cl(X)\) determines a character
\[
\chi_\alpha:G\longrightarrow\Gm,
\qquad
\chi_\alpha(g)=g(\alpha).
\]
Use this character to let \(G\) act on one extra affine coordinate:
\[
g\cdot(z,t)
=
\bigl(g\cdot z,\chi_\alpha(g)t\bigr),
\qquad
(z,t)\in\widehat X\times\AAA^1.
\]
The goal is to understand the quotient
\[
L_\alpha
=
(\widehat X\times\AAA^1)/G
\]
and to show that it is the total space of the line bundle represented
by \(\alpha\).

\begin{figure}[ht]
\centering
\begin{tikzcd}[column sep=large,row sep=large]
\widehat X\times\AAA^1
  \arrow[r,"{/G}"]
  \arrow[d,"\operatorname{pr}_1"']
&
L_\alpha
  \arrow[d]
\\
\widehat X
  \arrow[r,"\pi"']
&
X
\end{tikzcd}
\caption{A character of the Cox group tells the extra coordinate how
to rescale. Quotienting the trivial line over \(\widehat X\) then
produces a line bundle over \(X\).}
\label{fig:line-bundle-extra-cox-coordinate}
\end{figure}

\begin{enumerate}[label=\textbf{Hint \arabic*.},wide=0pt]

\item \textbf{The quotient still lies over \(X\).}

Show that the projection
\[
\operatorname{pr}_1:\widehat X\times\AAA^1\longrightarrow\widehat X
\]
is \(G\)-equivariant. Deduce a map
\[
p_\alpha:L_\alpha\longrightarrow X.
\]
Explain why the class of \((z,0)\) gives a zero section. Why should the
fiber over \(\pi(z)\) be a one-dimensional vector space rather than a
copy of \(\AAA^1\) with no preferred linear structure?

\item \textbf{Calibrate the sign on \(\PP^1\).}

For \(X=\PP^1\), we have
\[
\widehat X=\AAA^2\setminus\{0\},
\qquad
G=\Gm.
\]
For \(d\in\ZZ\), consider
\[
\lambda\cdot(x_0,x_1,t)
=
(\lambda x_0,\lambda x_1,\lambda^d t).
\]
On the chart \(x_i\neq0\), show that
\[
\tau_i=\frac{t}{x_i^d}
\]
is invariant. On the overlap, compute
\[
\tau_1
=
\left(\frac{x_0}{x_1}\right)^d\tau_0.
\]
Recognize the quotient as the total space of
\[
\mathcal O_{\PP^1}(d).
\]

\item \textbf{Build a local homogeneous unit.}

Choose a torus-invariant divisor
\[
D=\sum_\rho a_\rho D_\rho
\]
with \([D]=\alpha\). Since \(X\) is smooth, \(D\) is Cartier. For a
maximal cone \(\sigma\), let \(m_\sigma\in M\) be its Cartier datum,
so
\[
\langle m_\sigma,u_\rho\rangle=-a_\rho
\qquad
(\rho\in\sigma(1)).
\]
On
\[
\widehat U_\sigma=D(x^{\widehat\sigma}),
\]
define the Laurent monomial
\[
q_\sigma
=
\prod_\rho
x_\rho^{a_\rho+\langle m_\sigma,u_\rho\rangle}.
\]
Show that:
\begin{enumerate}[label=(\alph*),wide=0pt]
\item \(q_\sigma\) has Cox degree \(\alpha\);
\item its exponent is zero for every \(\rho\in\sigma(1)\);
\item it is therefore an invertible function on
\(\widehat U_\sigma\);
\item it transforms by
\[
q_\sigma(g\cdot z)=\chi_\alpha(g)q_\sigma(z).
\]
\end{enumerate}

\item \textbf{Trivialize the quotient locally.}

On \(\widehat U_\sigma\times\AAA^1\), set
\[
\tau_\sigma=\frac{t}{q_\sigma}.
\]
Show that \(\tau_\sigma\) is \(G\)-invariant. Use it to construct an
isomorphism
\[
p_\alpha^{-1}(U_\sigma)
\simeq
U_\sigma\times\AAA^1.
\]
Conclude that \(p_\alpha:L_\alpha\to X\) is a line bundle.

\item \textbf{Read the transition functions.}

On \(U_\sigma\cap U_\tau\), show that
\[
\tau_\tau
=
\frac{q_\sigma}{q_\tau}\tau_\sigma.
\]
Compute
\[
\frac{q_\sigma}{q_\tau}
=
\prod_\rho
x_\rho^{\langle m_\sigma-m_\tau,u_\rho\rangle}.
\]
Use \cref{thm:degree-zero-localization-affine-chart} to identify this
degree-zero Laurent monomial with the torus character
\[
\character{m_\sigma-m_\tau}.
\]
Compare with the transition functions of \(\mathcal O_X(D)\), and
deduce
\[
\boxed{L_\alpha\simeq\mathcal O_X(D).}
\]

\item \textbf{Why only the grading matters.}

Let
\[
a=(a_\rho)_\rho,
\qquad
b=(b_\rho)_\rho
\]
be two exponent vectors in \(\ZZ^{\Sigma(1)}\) with the same
class-group degree. Show that
\[
a-b
=
\bigl(\langle m,u_\rho\rangle\bigr)_\rho
\]
for some \(m\in M\).

The vectors \(a,b\) define characters of the large coordinate torus:
\[
\widetilde\chi_a((t_\rho)_\rho)
=
\prod_\rho t_\rho^{a_\rho},
\qquad
\widetilde\chi_b((t_\rho)_\rho)
=
\prod_\rho t_\rho^{b_\rho}.
\]
Show that their ratio factors through
\[
(\Gm)^{\Sigma(1)}\longrightarrow T_N.
\]
Since \(G\) is the kernel of this map, conclude that
\[
\widetilde\chi_a|_G
=
\widetilde\chi_b|_G.
\]
Thus two choices with the same Cox grading give the same action on the
extra direction, hence isomorphic line bundles.

Interpret this as the familiar statement
\[
D\sim D'
\quad\Longrightarrow\quad
\mathcal O_X(D)\simeq\mathcal O_X(D').
\]

\item \textbf{Tensor products add weights.}

Construct natural isomorphisms
\[
L_\alpha\otimes L_\beta
\simeq
L_{\alpha+\beta},
\qquad
L_\alpha^\vee
\simeq
L_{-\alpha}.
\]
Explain why this is already visible from the action on the extra
coordinate.

\item \textbf{Homogeneous polynomials become sections.}

Let \(f\in(S_\Sigma)_\alpha\). Then
\[
f(g\cdot z)=\chi_\alpha(g)f(z).
\]
Show that
\[
\pi(z)
\longmapsto
[z,f(z)]
\]
defines a global section of \(L_\alpha\). Use
\cref{thm:cox-piece-global-sections} and the identification from Hint
5 to recover
\[
(S_\Sigma)_\alpha
\simeq
H^0(X,L_\alpha).
\]

\item \textbf{The smooth conclusion: \(\Cl=\Pic\).}

The construction gives a homomorphism
\[
\Phi:\Cl(X)\longrightarrow\Pic(X),
\qquad
\alpha\longmapsto[L_\alpha].
\]
Use the preceding hints to show that \(\Phi\) sends \([D]\) to
\([\mathcal O_X(D)]\). Explain why smoothness implies that every Weil
divisor is Cartier, and conclude that
\[
\boxed{
\Cl(X)\simeq\Pic(X).
}
\]
Where would this argument fail for a singular toric variety?

\item \textbf{Return to the running blowup.}

For
\[
X=\Blp{\PP^2}{p},
\qquad
\Cl(X)=\ZZ H\oplus\ZZ E,
\]
write \(g=(\lambda,\mu)\in(\Gm)^2\). The class
\[
\alpha=dH-rE
\]
gives the character
\[
\chi_\alpha(\lambda,\mu)=\lambda^d\mu^{-r}.
\]
Write the full action on
\[
\widehat X\times\AAA^1
\]
and describe the quotient line bundle. Explain how a Cox monomial of
degree \(dH-rE\) becomes a section of this bundle. Compare your answer
with the truncated polygon \(P_{dH-rE}\) from Chapter~5.

\end{enumerate}

\begin{tcolorbox}[
  colback=PredictionTeal!5,
  colframe=PredictionTeal,
  title={The lesson of the construction},
  fonttitle=\bfseries
]
The Cox quotient does more than recover the base variety. Once
\(\widehat X\to X\) is known, every character of the rescaling group
tells one additional coordinate how to transform. The resulting
associated line bundle depends only on the character---equivalently,
only on the divisor class. For a smooth toric variety, all divisor
classes are Cartier, so this construction realizes the entire Picard
group:
\[
\Cl(X)=\Pic(X)=\Hom(G,\Gm).
\]
\end{tcolorbox}

\noindent
\textit{Reference.}
See Cox--Little--Schenck, Lemma~5.1.1, Proposition~4.2.6, and
\S14.1.

\clearpage
\section*{Notes and references}
\addcontentsline{toc}{section}{Notes and references}
\markboth{Notes and references}{Notes and references}

The Cox quotient construction is developed in Cox--Little--Schenck,
\S5.1. The class-group grading and homogeneous coordinate ring are
treated in \S5.2, homogeneous ideals in \S5.3, and graded modules and
sheaves in \S5.4. The local identity
\[
\left((S_\Sigma)_{x^{\widehat\sigma}}\right)_0
\simeq
\Bbbk[\sigma^\vee\cap M]
\]
is the algebraic core of the quotient construction.

The relation between smoothness, factoriality, and
\(\Pic(X_\Sigma)=\Cl(X_\Sigma)\) is developed in Chapter~4. The
challenge exercise recasts that equality through associated bundles:
for a free Cox quotient, characters of \(G\) produce line bundles on
the quotient. This is also the basic linearization idea used in toric
geometric invariant theory; see Cox--Little--Schenck, \S14.1.

The phrase \emph{universal torsor} is often used for the free Cox
quotient
\[
\widehat X_\Sigma\longrightarrow X_\Sigma
\]
of a smooth toric variety. No torsor formalism was needed here: the
concrete content is the familiar associated-bundle construction
\[
(\widehat X_\Sigma\times\AAA^1)/G.
\]

Later chapters can use Cox coordinates not only to describe the
ambient toric variety, but also to write and study subvarieties inside
it by means of homogeneous equations.
